\documentclass[11pt,a4paper]{article}
\usepackage[margin=27mm,headheight=14pt]{geometry}
\usepackage[T1]{fontenc}
\usepackage{lmodern,amsmath,amssymb,amsthm,mathtools,mathrsfs}
\usepackage{microtype,booktabs,longtable,array,enumitem,xcolor,graphicx}
\usepackage{xurl}
\usepackage[colorlinks=true,allcolors=blue!45!black]{hyperref}
\usepackage{bookmark}
\setlength{\emergencystretch}{2em}
\setlength{\parskip}{3pt}
\setlist{itemsep=2pt,topsep=4pt}
\allowdisplaybreaks[2]
\newtheorem{theorem}{Theorem}[section]
\newtheorem{proposition}[theorem]{Proposition}
\newtheorem{lemma}[theorem]{Lemma}
\newtheorem{corollary}[theorem]{Corollary}
\newtheorem{counterexample}[theorem]{Counterexample}
\theoremstyle{definition}
\newtheorem{definition}[theorem]{Definition}
\newtheorem{remark}[theorem]{Remark}
\newcommand{\Sh}{\operatorname{Sh}}
\newcommand{\Hom}{\operatorname{Hom}}
\newcommand{\Map}{\operatorname{Map}}
\newcommand{\Pro}{\operatorname{Pro}}
\newcommand{\colim}{\operatorname*{colim}}
\newcommand{\limm}{\operatorname*{lim}}
\newcommand{\id}{\operatorname{id}}
\newcommand{\Spec}{\operatorname{Spec}}
\newcommand{\im}{\operatorname{im}}
\newcommand{\cl}{\operatorname{cl}}
\newcommand{\R}{\mathbb R}
\newcommand{\C}{\mathbb C}
\newcommand{\Z}{\mathbb Z}
\newcommand{\Q}{\mathbb Q}
\newcommand{\F}{\mathbb F}
\newcommand{\recordurl}{https://github.com/KokunoYumeto/mathematics-workbenches}
\newcommand{\workbenchurl}{https://github.com/KokunoYumeto/homotopy-and-categories-workbench}
\title{Grothendieck's interval supports and oriented maquettes,\\ after Jen\v ca and Bj\"orner}
\author{}
\date{23 November 2026}
\begin{document}
\maketitle
\begin{abstract}
Grothendieck's manuscripts on a geometry of forms associate support lattices to disjunction relations and seek geometric interpretations of figures and oriented incidence data. Starting from Jen\v ca's theorem for finite bounded posets, we characterize orthomodular interval supports for every finite poset, including its points, and give a Boolean criterion without a finiteness assumption. We separate finite figure-complement identities from global orthomodularity, give their maximal-clique criterion, and describe complete support maps by saturated correspondences. We identify the precise distributivity condition under which finite figures are configurations of binary-conflict event structures. Exact representation theorems compare carrier axioms with internal factorization and admissible set figures, including the indispensable principal-subfigure closure. With nonempty open lieux and witness absorption, full figure atlases admit subdivision extensions, while uniqueness of a deployment diagram does not imply uniqueness of its realization. Bj\"orner's characterization of regular cell posets, together with a signed-boundary argument, gives the corresponding oriented characterization in all dimensions. A torus boundary shows why connectedness and signed diamonds cannot replace spherical lower intervals. Simplicial recognition requires a global intersection condition, while the order-complex map to an Alexandrov poset is generally only a weak equivalence. We also construct a definable cell atelier over every o-minimal real closed field, prove its finite common-refinement and coarsest boundary-extension properties, and compute its complete support algebra. The mathematics and text were produced by GPT-6 Astra (ChatGPT Pro mode), with checks recorded by scope; no human has checked the work.
\end{abstract}

\section{Supports, intervals and the main statements}
This section specifies the structures being classified. An \emph{orthoset} is a set $E$ with a symmetric irreflexive relation $\perp$. For $A\subset E$ put
\[
 N(A)=\{x:(\forall a\in A)\ x\perp a\},\qquad
 \cl A=N^2(A),\qquad\Sigma(E)=\{A:\cl A=A\}.
\]
The complete support lattice has intersections as meets, closure of unions as joins, and $A\mapsto N(A)$ as orthocomplement. Its least and greatest elements are $\varnothing,E$. These are the polar constructions used in cote~156-9. On archive sheet~52 (PDF p.~53, manuscript p.~49), Grothendieck writes that to an ``atelier spatial'' there ``est associ\'e un Spatos'' \cite{Groth}. The support algebra is therefore part of the proposed passage from elementary geometric objects to a global spatial structure. We ask when this algebra has the orthomodular or Boolean laws, and which additional geometric hypotheses are needed to realize incidence data by cells.

Hua's account of cote~156-4 gives a seven-element compatibility poset whose complete support lattice is a nondistributive hexagon \cite[Counterexample~79.1, pp.~369--370]{Hua}. This already disproves a general Boolean conclusion under those compatibility assumptions. The interval classification below concerns the explicitly specified disjunction relation, and the definable atelier later in the paper has a different, primitive disjunction of open cells.

For a poset $L$, its \emph{interval magasin with points} is
\[
 M(L)=\{P_x:x\in L\}\amalg\{I_{ab}:a<b\}.
\]
The orthogonality relation is specified completely by
\begin{align*}
 P_x\perp P_y&\iff x<y\text{ or }y<x,\\
 P_x\perp I_{ab}&\iff x\le a\text{ or }b\le x,\\
 I_{ab}\perp I_{cd}&\iff b\le c\text{ or }d\le a.
\end{align*}
A \emph{conditional lattice} is a poset in which every upper-bounded pair has a join and every lower-bounded pair has a meet.

\begin{theorem}[The interval classification]\label{thm:interval}
For every finite poset $L$, the ortholattice $\Sigma(M(L))$ is orthomodular if and only if $L$ is a conditional lattice. For an arbitrary poset $L$, the ortholattice $\Sigma(M(L))$ is Boolean if and only if the nonisolated elements of $L$ form a chain. Here an element is nonisolated if it is comparable with a distinct element.
\end{theorem}

Jen\v ca's bounded theorem concerns proper quotients $[a<b]$ of a finite bounded poset \cite[Theorem~4.10, preprint p.~10]{Jenca}. We use that theorem only with its finiteness and boundedness hypotheses. The additional argument splits every point into two adjacent points, compares the resulting quotient logic with $M(L)$, and treats the new endpoints created by bounding an unbounded poset.

\begin{theorem}[Boolean and atom tests]\label{thm:tests}
For any orthoset $E$, its support lattice is Boolean exactly when
\[
 x\not\perp y\Longrightarrow
 (\exists z\in E)\ N(x)\cup N(y)\subseteq N(z).
\]
The nonzero support $\cl\{x\}$ is an atom exactly when
\[
 N(x)\subseteq N(y)\Longrightarrow N(x)=N(y)
 \quad\text{for every }y\in E.
\]
These assertions allow infinite orthosets.
\end{theorem}
The Boolean equivalence is Jen\v ca's general compatibility criterion \cite[Theorem~3.2, p.~5]{Jenca2026}; we include its proof to fix the conventions. His later correction concerns a different assertion, the forbidden-configuration characterization for incomparability orthosets \cite[Theorem~1, pp.~2--4]{JencaCorrection}. The correction adds an $X$ configuration to the forbidden weak $N$; it does not alter the general criterion or his orthomodularity theorem.

Call a pairwise orthogonal subset $F$ of an orthoset an
\emph{admissible figure} when
\[
 \cl F\cap N(A)=\cl(F\setminus A)\qquad(A\subseteq F).
\]
The quantifier over figures is important: their supports need not
represent every orthogonal family in the complete lattice.

\begin{theorem}[Figure complements and finite flags]\label{thm:flags}
For any orthoset $E$, the following are equivalent: every pairwise
orthogonal subset is admissible; and, for every maximal pairwise
orthogonal subset $D$ and distinct $u,v\in E$,
\[
 D\subseteq N(u)\cup N(v)\quad\Longrightarrow\quad u\perp v.
\]
The equivalence includes infinite subsets, with the axiom of choice.
For every poset $L$, each finite pairwise orthogonal subset of $M(L)$
is admissible. This finite assertion does not imply orthomodularity
of $\Sigma(M(L))$, and it does not extend to all infinite figures.
\end{theorem}
For finite graphs the displayed condition is Sumner's definition
\cite[p.~492]{Sumner}. It is the normality condition of
Paseka--Vetterlein in their finite-rank setting
\cite[Definition~2.4 and Lemma~2.5, pp.~5--6]{PV}.
It is weaker than the basis criterion for an orthomodular support
lattice. Section~\ref{sec:figure-comparisons} gives explicit finite
and infinite witnesses.

For subsets $A,B$ of an orthoset, $A\perp B$ means $A\subseteq N(B)$.
The object and arrow conventions also matter. A complete
ortholattice is a complete lattice with an antitone involution $C$
satisfying $x\wedge Cx=0$ and $x\vee Cx=1$. For orthosets $E,F$,
a \emph{saturated correspondence} is a map
$\alpha:E\to\Sigma(F)$ such that, for every $A\subseteq E$,
\[
 x\in\cl_E A\quad\Longrightarrow\quad
 \alpha(x)\subseteq\bigvee_{a\in A}\alpha(a).
\]

\begin{theorem}[The complete support correspondence]\label{thm:correspondence}
Every complete ortholattice $K$ is represented by the orthoset
$K\setminus\{0\}$, with $x\perp y$ if and only if $x\le Cy$:
\[
 K\xrightarrow{\sim}\Sigma(K\setminus\{0\}),\qquad
 s\longmapsto\{x\ne0:x\le s\}.
\]
For every pair of orthosets $E,F$, saturated correspondences are
in bijection with maps $f:\Sigma(E)\to\Sigma(F)$ preserving all
joins, including the empty join. The inverse formulas are
\[
 f(S)=\bigvee_{x\in S}\alpha(x),\qquad
 \alpha(x)=f(\cl_E\{x\}).
\]
The identity is $x\mapsto\cl_E\{x\}$ and composition is
$(\beta\alpha)(x)=\bigvee_{y\in\alpha(x)}\beta(y)$.
Consequently these two categories are equivalent. If $g$ is the
order right adjoint of $f$, the associated dagger is
$f^\dagger=C_EgC_F$; it is characterized by
\[
 f(S)\perp T\quad\Longleftrightarrow\quad S\perp f^\dagger(T).
\]
Ordinary orthogonality-preserving functions are not the arrows in
this equivalence.
\end{theorem}
The object axioms are exactly the Spat axioms on archive
sheets~50--52 of cote~156-9 \cite{Groth}. We specify the arrow
condition explicitly so that this statement does not depend on an
unspecified meaning of morphism.


For finite figures we use a poset $\mathcal F$ with zero, bounded suprema, and its nonzero completely join-irreducible elements $M$. Write $\delta(F)=\{m\in M:m\le F\}$. Assume every $\delta(F)$ is finite, $F=\bigvee\delta(F)$, and two figures have a common upper bound exactly when each pair of their deployed irreducibles has a common upper bound.

\begin{theorem}[The event-structure class]\label{thm:events}
Under these hypotheses, the following are equivalent: every principal ideal of $\mathcal F$ is distributive; every $m\in M$ is prime for bounded joins; and deployment identifies $\mathcal F$ with the finite downward-closed pairwise compatible subsets of $M$. Equivalently, these are exactly the finite configuration posets of binary-conflict prime event structures with finite causes. Conflict is failure to have a common upper bound, and causality is the order on $M$.
\end{theorem}

The oriented result uses a set $M$ with a free reversal involution
$m\mapsto\bar m$ and an immediate incidence relation satisfying
$n\prec m$ if and only if $\bar n\prec\bar m$. Its projected
relation on the reversal orbits generates a poset $P$. Distinguish
one positive lift at each minimal orbit, to be the positive
zero-cell orientation. Choose one lift $x^+$ of every other orbit.
When a cover $y<x$ has a unique incident lower orientation for
$x^+$, define $\epsilon(y,x)=+1$ if $y^+\prec x^+$ and
$\epsilon(y,x)=-1$ if $\overline{y^+}\prec x^+$.
For a poset $Q$, $\Delta(Q)$ is the simplicial complex of its
finite nonempty chains. Its realization has the weak topology.
We use $S^{-1}=\varnothing$.

\begin{theorem}[Oriented regular cell structures]\label{thm:cw}
The oriented data arise from a regular CW complex with the specified
positive vertex orientations if and only if:
\begin{enumerate}
\item every principal ideal of $P$ is finite;
\item immediate incidences project exactly to covers and, for each
cover and each upper orientation, exactly one lower orientation
is incident to it;
\item for each $x\in P$ there is an integer $d(x)\ge0$ such that
$|\Delta(P_{<x})|$ is homeomorphic to $S^{d(x)-1}$;
\item the two incidence-sign products around every rank-two
 diamond are opposite, and every edge has opposite coefficients
at its two vertices relative to their distinguished orientations.
\end{enumerate}
The spherical condition makes $P$ graded with dimension rank $d$:
every maximal chain from a minimal element to $x$ has $d(x)$
strict steps, every rank-two interval has exactly two intermediate
elements, and every edge has exactly two vertices.
The conclusion allows infinitely many cells, unbounded dimension
and failure of local finiteness; the realization has the CW weak
topology. The empty data give the empty CW complex.
\end{theorem}

The unoriented characterization is Bj\"orner's theorem \cite[Definition~2.1 and Proposition~3.1]{Bjorner}. The additional point is that the prescribed signs are realized, not merely that some orientations can be chosen. No claim of a category equivalence with unspecified morphisms is made.

For a poset $P$, write $P^0$ for the poset obtained by adjoining a new
minimum $0$. The lower Alexandrov topology on $P$ has the principal
lower sets $P_{\le x}$ as a basis.

\begin{theorem}[Simplicial recognition and Alexandrov comparison]
\label{thm:simplicial-comparison}
A poset $P$ is the nonempty face poset of an abstract simplicial
complex if and only if every $[0,x]\subseteq P^0$ is a finite Boolean
lattice and $P^0$ has binary meets. Equivalently, under the same
finite Boolean interval hypothesis, the sets of atoms below distinct
elements of $P$ are distinct.

For every poset $L$, without a finiteness or local finiteness
assumption, the map
\[
 \mu_L:|\Delta(L)|\longrightarrow L,\qquad
 \mu_L(p)=\min\operatorname{supp}(p),
\]
to its lower Alexandrov topology is a weak homotopy equivalence.
It need not be a homotopy equivalence or a homeomorphism.
\end{theorem}
The first equivalence is Bj\"orner's Boolean-type criterion
\cite[Section~2.3, pp.~8--9]{Bjorner}; the second assertion is
McCord's theorem \cite[Theorem~2, p.~466; Section~5, p.~469]{McCord}.
Section~\ref{sec:simplicial-comparison} proves the recognition
criterion, checks McCord's hypotheses, and distinguishes the two
flag constructions appearing in the manuscripts.


\paragraph{The definable cell model.}
Fix an o-minimal expansion of a real closed field $R$ and a definable set $X\subseteq R^n$. Parameters from $R$ are allowed. All closures in this section are taken in the fixed space $X$. In particular, neither $X$ nor the support of a figure need be closed in $R^n$ or bounded.

\begin{definition}
A \emph{generalized cell} is a nonempty definable set definably homeomorphic to $(0,1)^d$, with a singleton for $d=0$. A \emph{cell figure} is a finite partition of a relatively closed definable $S\subseteq X$ into generalized cells, subject to the frontier condition
\[
 D\cap\overline C^{\,X}\ne\varnothing\quad\Longrightarrow\quad D\subseteq\overline C^{\,X}.
\]
The empty figure is allowed. Write $|F|=S$. A figure can equivalently be specified by its finite family of closed strata $A=\overline C^{\,X}$, ordered by inclusion. Its open stratum is
\[
 C_A=A\setminus\bigcup_{B\subsetneq A,\ B\in F}B.
\]
An \emph{elementary figure}, or multistrate, is the family of all closed strata below one $A$. It remembers its entire boundary decomposition; its support is $A$ and its top cell is $m^\circ=C_A$.
\end{definition}
The frontier-dimension theorem implies that every proper boundary cell has lower dimension \cite[Theorem~3.22, pp.~34--35]{Coste}. Thus the closed strata of different cells are distinct, the displayed reconstruction is valid, and $C_A$ is relatively open and dense in $A$. We impose no closed-ball condition on $A$.

There are two orders. Incidence $G\le F$ means downward-closed inclusion of closed-stratum families. Refinement $G\ll F$ means that each open cell of $G$ is contained in an open cell of $F$. A subdivision $G\preceq F$ is a refinement with $|G|=|F|$. Compatibility means existence of a common upper figure for incidence. On elementary figures define primitive disjunction by
\[
 m\perp^\circ n\quad\Longleftrightarrow\quad m^\circ\cap n^\circ=\varnothing.
\]



For $G\ll F$, write $c_{G,F}(B)$ for the unique target closed stratum whose open cell contains $C_B$. For $H\le F$ put $G|H=c_{G,F}^{-1}(H)$. The boundary of an elementary figure is its subfigure of proper closed strata. The carrier calculus establishing these definitions is proved in Section~\ref{sec:cellatelier}.

\begin{theorem}[Coarsest extension of a boundary subdivision]\label{thm:extension}
Let $G\le F$ and let $G'\preceq G$. There is a coarsest subdivision $E\preceq F$ inducing $G'$ over $G$. Its closed-stratum family is
\[
 E=G'\cup(F\setminus G).
\]
The open cells outside $|G|$ remain the old open cells. Every subdivision $H\preceq F$ with $H|G=G'$ satisfies $H\ll E$. For a new stratum $B\in G'$ and an old stratum $A\in F\setminus G$,
\[
 B\subseteq A\quad\Longleftrightarrow\quad c_{G',G}(B)\subseteq A.
\]
In particular, subdividing the boundary of an elementary figure extends to an elementary figure with the same top cell.
\end{theorem}

\begin{theorem}[Finite common refinements]\label{thm:moderation}
Every relatively closed definable $S\subseteq X$ admits a cell figure. Any finite collection of cell figures on the same support has a common subdivision. In particular, any finite collection of subdivisions of one figure has a common subdivision.
\end{theorem}

\begin{theorem}[Complete supports and spatial detection]\label{thm:cellpowerset}
Let $M_X$ be the elementary figures with primitive disjunction $\perp^\circ$. For an arbitrary subset $Q\subseteq X$, put
\[
 S_Q=\{m\in M_X:m^\circ\subseteq Q\}.
\]
Then $Q\mapsto S_Q$ is an isomorphism from the complete Boolean algebra $\mathcal P(X)$ onto $\Sigma(M_X)$. No definability assumption on $Q$ is imposed. An incidence subfigure has the same total support whether computed from its open strata or its elementary figures. Moreover $|F|\subseteq|G|$ if and only if every singleton figure refining $F$ also refines $G$.
\end{theorem}

\paragraph{The abstract carrier axioms.}
We also give exact representations of the incidence and refinement
data themselves. A \emph{prime carrier atelier} is a nonempty poset
$\mathscr F$ with an additional refinement order $\ll$, satisfying
the following conditions. The least element is denoted $0$.
\begin{enumerate}
\item Every bounded family, including the empty family, has a join.
The nonzero completely join-prime elements form a set $M$, and
every $F\in\mathscr F$ is the join of
$\widetilde F=\{x\in M:x\leq F\}$.
Primality here means that $x\leq\bigvee_iF_i$ for a bounded
family implies $x\leq F_i$ for some $i$.
\item Two figures $F,G$ have a common upper bound if and only if
every $x\in\widetilde F$ and $y\in\widetilde G$ have one.
Such figures are called compatible.
\item Refinement is a partial order containing incidence, and
\[
 F\ll G\quad\Longleftrightarrow\quad
 (\forall x\in\widetilde F)(\exists y\in\widetilde G)\ x\ll y.
\]
\item If $x\in M$ refines $F$, then
$\{y\in\widetilde F:x\ll y\}$ has a least element, the carrier
$c_F(x)$.
\item Compatible refinement is incidence: $F\ll G$ and
compatibility of $F,G$ imply $F\leq G$.
\item The internal relation $x\rho y$, defined by $x\ll y$
and $c_y(x)=y$, is transitive on $M$.
\end{enumerate}
All sets and figures may be infinite. Isomorphisms preserve both
orders. These are explicit hypotheses; no chronological axiom
label is needed to interpret them.

\paragraph{Two strong presentations on the strata.}
A \emph{strong carrier system} consists of partial orders
$\leq\subseteq\ll$ on a set $M$, least carriers
$c(x,y)=\min\{z\leq y:x\ll z\}$ whenever $x\ll y$,
and a symmetric irreflexive relation $\delta$. Put
$x\rho y$ when $c(x,y)=y$. Require $\delta$ to be hereditary
under $\rho$ in both variables and require
\begin{equation}\label{eq:principal-disjoint}
 y,z\leq x,\quad y\ne z\quad\Longrightarrow\quad y\delta z.
\end{equation}
A \emph{strong factorization system} instead has partial orders
$\leq,\rho$, the same disjunction conditions, at most one $z$
with $x\rho z\leq y$ for any fixed $x,y$, and exchange
\begin{equation}\label{eq:carrier-exchange}
 x'\leq x\rho y\quad\Longrightarrow\quad
 \exists y'\leq y\ (x'\rho y').
\end{equation}

\begin{theorem}[The strong relational comparison]
\label{thm:strong-relations}
The two presentations are equivalent by mutually inverse definitions
on the same set. The inverse refinement formula is
\[
 x\ll y\quad\Longleftrightarrow\quad
 \exists z\ (x\rho z\leq y).
\]
For either presentation, let $\mathscr H$ be a family of
downward-closed $\delta$-cliques containing $\varnothing$ and every principal downset,
closed under downward-closed subfigures and under finite unions
that remain cliques. With incidence inclusion and refinement tested
elementwise, $\mathscr H$ is a prime carrier atelier.

Write $P|G=\{p\in P:c_F(p)\in G\}$ when $P\ll F$ and
$G\subseteq F$. If $F=\bigcup_{i=1}^nF_i$ and
$P_i\ll F_i$ satisfy
\[
 P_i|(F_i\cap F_j)=P_j|(F_i\cap F_j)\quad\text{for all }i,j,
\]
there is a unique $P\ll F$ with $P|F_i=P_i$, namely
$P=\bigcup_iP_i$. For the full family of clique downsets,
the same conclusion holds for arbitrary set-indexed covers,
and every directed incidence union is admitted.
\end{theorem}



\paragraph{Admissible set figures.}
A family $\Phi$ of nonempty subsets of $E$ is \emph{admissible}
if every point of $\bigcup\Phi$ belongs to a least member and
every member occurs as the least member of some point.
Its nonempty cells are
\[
 C_A^\Phi=A\setminus\bigcup_{B\in\Phi,\,B\subsetneq A}B.
\]
Write $\operatorname{Fig}(E)$ for these families, including the
empty one. Set-figure refinement means that every source cell
is contained in a target cell. Put
$\Phi_A=\{B\in\Phi:B\subseteq A\}$.

The representation is already formulated in Grothendieck's
\emph{Th\'eor\`eme-Scholie}, archive sheets~32--33 of part~6
\cite[PDF pp.~33--34]{Groth6}. The principal-membership clause
appears expressly in the preceding formulation on archive sheet~31
(PDF p.~32). We state the hypotheses uniformly and prove the
equivalence for isomorphisms.

\begin{theorem}[The set-atlas representation]
\label{thm:set-atlas}
Prime carrier ateliers are equivalent, as a groupoid, to pairs
$(E,\mathscr H\subseteq\operatorname{Fig}(E))$ with these properties:
\begin{enumerate}
\item $\varnothing\in\mathscr H$, and every union bounded by a
member of $\mathscr H$ belongs to $\mathscr H$.
\item If $A\in\Phi\cap\Psi$, then
$\Phi_A=\Psi_A$ and this common principal subfigure belongs
to $\mathscr H$. Denote it by $\widehat A$.
\item For occurring strata $A\subseteq B$, one has
$\widehat A\ll\widehat B$ in $\operatorname{Fig}(E)$.
\item If $\widehat A\cup\widehat B\in\mathscr H$ for every
$A\in\Phi,B\in\Psi$, then $\Phi\cup\Psi\in\mathscr H$.
\item The global family
$\mathscr U=\bigcup_{\Phi\in\mathscr H}\Phi$ is admissible,
has support $E$, and separates its points.
\end{enumerate}
Morphisms of pairs are bijections of $E$ transporting the admitted
families. The canonical forward construction is
\[
 E=M,\qquad A_x=\{z\in M:z\ll x\},\qquad
 \Phi_F=\{A_x:x\leq F\}.
\]
The principal-membership clause in condition~2 is indispensable.
The conclusion includes the empty ground set and does not assert
an equivalence for a larger, unspecified class of morphisms.
\end{theorem}

\paragraph{Detection by open lieux.}
In either presentation a \emph{lieu} is a minimal refinement
stratum. Write $L$ for their set,
$U_x=\{p\in L:p\rho x\}$, and
$\ell(F)=\{p\in L:p\ll F\}$.
For figures, write $F\parallel G$ when they are compatible
and their incidence meet is zero.

\begin{theorem}[The local-detection equivalence]
\label{thm:local-relations}
Strong systems with admitted atlases as in
Theorem~\ref{thm:strong-relations} and satisfying
\[
 x\delta y\quad\Longleftrightarrow\quad
 (\forall p\in U_x)(\forall q\in U_y)\ p\delta q
\]
are equivalent, on objects and isomorphisms, to prime carrier
ateliers satisfying the following conditions for strata $x,y$:
\begin{enumerate}
\item $U_x\ne\varnothing$.
\item $x\parallel y$ if and only if every
$p\in\ell(x),q\in\ell(y)$ satisfies $p\parallel q$.
\item If $x,y$ are distinct and compatible, every
$p\in U_x,q\in U_y$ satisfies $p\parallel q$.
\item With $K=x\wedge y$, if every
$p\in\ell(x)\setminus\ell(K)$ and
$q\in\ell(y)\setminus\ell(K)$ satisfies $p\parallel q$,
then $x,y$ are compatible.
\end{enumerate}
The inverse primitive relation is defined by the same universal
open-lieu formula, using $p\parallel q$ as the graph on lieux.
No finiteness assumption is imposed, and the empty stratum set
is allowed.

More precisely, suppose only the incidence/refinement partial
orders and least carriers of a strong carrier system are given,
together with nonempty $U_x$, a symmetric irreflexive graph on
$L$, and principal disjointness~\eqref{eq:principal-disjoint}
tested by that graph on open lieux. The induced relation is
strong if and only if internal refinement is transitive.
It suffices to impose
$p\rho x\rho y\Longrightarrow p\rho y$ for $p\in L$;
this restricted condition forces full transitivity.
In the factorization presentation, nonempty $U_x$ and the
all-pairs principal condition are the corresponding exact inputs.
\end{theorem}

\paragraph{Subdivision and absorption.}
For an abstract figure put
$\mathcal N(F)=\{H:H\parallel F\}$, and define
$P\preceq F$ by $P\ll F$ and
$\mathcal N(P)=\mathcal N(F)$.
The \emph{witness-absorption condition} says that if $P\ll F$
and some $H$ is disjoint from $P$ but not from $F$, such a witness
can be chosen to refine $F$.

\begin{theorem}[Extension in the full figure atlas]
\label{thm:abstract-extension}
Take a strong system of Theorem~\ref{thm:strong-relations} with
the full family of downward-closed disjunction cliques as figures.
Assume that every open-lieu set $U_x$ is nonempty and that
witness absorption holds. Then $P\preceq F$ if and only if
$P$ is incidence-maximal among figures refining $F$.
For every incidence subfigure $G\leq F$ and every subdivision
$G'\preceq G$, there is a subdivision $F'\preceq F$ with
\[
 F'|G=G'.
\]
The existence proof uses the axiom of choice through Zorn's lemma.
There is no finiteness hypothesis. Uniqueness of the extension
does not follow, even if its ordered deployment and carrier map
are prescribed up to their fixed-piece isomorphisms.
\end{theorem}

\paragraph{The interval carrier data.}
The point-and-interval construction at the start of this section
also has the incidence and refinement structure of the manuscript
\cite[part~9, archive sheets~4--7, PDF pp.~5--8]{Groth}.
Incidence on $M(L)$ has identities and
$P_a,P_b\leq I_{ab}$. Internal refinement has identities and
\[
 P_x\rho I_{ab}\iff a<x<b,\qquad
 I_{cd}\rho I_{ab}\iff a\leq c<d\leq b,
\]
with no interval internally refining a point.
Write $\lambda(P_x)=\mu(P_x)=x$,
$\lambda(I_{ab})=a$, and $\mu(I_{ab})=b$.
The positional order is
$U<_{\mathrm p}V$ if $U\ne V$ and
$\mu(U)\leq\lambda(V)$.
A strong system is called \emph{faithful} when two strata are
disjoint exactly when they have no common internal refinement:
\[
 x\delta y\quad\Longleftrightarrow\quad
 \nexists z\ (z\rho x\text{ and }z\rho y).
\]

\begin{theorem}[The complete interval structure]
\label{thm:interval-system}
For every poset $L$, these incidence and internal-refinement data,
with the orthogonality of $M(L)$ as primitive disjunction, form
a strong system. Its overall refinement is
\[
 \begin{aligned}
 P_x\ll P_y&\iff x=y,\\
 P_x\ll I_{ab}&\iff a\leq x\leq b,\\
 I_{cd}\ll I_{ab}&\iff a\leq c<d\leq b,
 \end{aligned}
\]
and no interval refines a point. It is faithful if and only if
$L$ is a chain, including the empty and singleton cases.
Its full figures are exactly the positional chains containing
the endpoints of all their intervals; its finite atlas consists
of the finite such chains. No local finiteness or density is
assumed.

Its lieux are precisely the point symbols, and
\[
 U_{P_x}=\{P_x\},\qquad U_{I_{ab}}=\{P_x:a<x<b\}.
\]
Thus all open-lieu sets are nonempty exactly when every strict
comparable pair in $L$ has a strict intermediate point.
\end{theorem}

\paragraph{Spatial completion.}
The next construction uses the manuscript's inclusion-carrier
condition \cite[part~9, archive sheet~45, PDF p.~46]{Groth}
and its target of nonzero orthogonal labels
\cite[archive sheet~47, PDF p.~48]{Groth}.
For an orthoset $(E,\perp)$ with a partial incidence order $\leq$,
put $S_x=\cl\{x\}$ and label the principal ideal $[x]$ by
$z\mapsto S_z$, retaining its incidence order. Write $D_x$ for
this ordered labelled deployment. An \emph{inclusion carrier}
$h:D_x\to D_y$ is order-preserving and satisfies $S\subseteq h(S)$
for every label. It is internal if $h(S_x)=S_y$.

\begin{theorem}[Completion in the spatial target]\label{thm:spatial-completion}
Assume that, in each $D_x$, the labels are nonzero, distinct and
pairwise orthogonal, and that the ordered labelled deployment
$D_x$ determines $x$ injectively. Declare $x\rho y$ exactly when
an internal inclusion carrier $D_x\to D_y$ exists. Then this is
a strong system with the given incidence and orthogonality.
Its overall refinement is exactly existence of an inclusion
carrier without the internal condition. Each carrier is unique
and factors uniquely through the principal ideal of the image
of its top label. This is the unique completion whose internal
arrows are all the internal inclusion carriers in this specified
spatial target. No finiteness assumption is needed.

Pairwise orthogonality cannot be replaced by mere disjointness
of the label subsets. The latter permits five-point data for
which the principal-disjunction axiom is impossible.
\end{theorem}

\paragraph{Relative and absolute criteria.}
The following comparison is for a prime carrier atelier, so it allows
restricted atlases. Write $F\bowtie G$ for compatibility and
$F\parallel G$ for compatibility together with $F\cap G=0$.
Thus zero is disjoint from every figure, including itself.
Let $\ell(F)=\{p\in L:p\ll F\}$ and $U_x=\{p:p\rho x\}$.
For compatible parents $F,G$, put $K=F\cap G$ and take
$P\ll F,Q\ll G$. Define the relative conditions
\[
\begin{aligned}
 \mathrm C &: P|K\bowtie Q|K\ \Longrightarrow\ P\bowtie Q,\\
 \mathrm D &: P|K\parallel Q|K\ \Longrightarrow\ P\parallel Q,\\
 \mathrm D_0 &: F\parallel G\ \Longrightarrow\ P\parallel Q.
\end{aligned}
\]
The three lieu conditions, for elementary $x,y$, are
\[
\begin{aligned}
 \mathrm L1 &: U_x\ne\varnothing,\\
 \mathrm L2 &: x\parallel y\ \Longleftrightarrow\
  (\forall p\in\ell(x),q\in\ell(y))\ p\parallel q,\\
 \mathrm L3 &: x\bowtie y,\ x\ne y\ \Longrightarrow\
  (\forall p\in U_x,q\in U_y)\ p\parallel q.
\end{aligned}
\]
For arbitrary elementary $x,y$, put $K=x\cap y$. The absolute
condition $\mathrm{CL}$ is
\[
 (\forall p\in\ell(x)\setminus\ell(K),\ q\in\ell(y)\setminus\ell(K))
       \ p\parallel q\quad\Longrightarrow\quad x\bowtie y.
\]
The condition $\mathrm{CD2}$ replaces this antecedent by
\[
 (\forall P\ll x,Q\ll y)\quad
 [P|K=Q|K=0\ \Longrightarrow\ P\parallel Q].
\]
The polyhedral sufficient versions $\mathrm{CL}_{\rm p}$ and
$\mathrm{CD2}_{\rm p}$ impose, in addition, that $K$ be empty or
an elementary principal figure. They do not assert that every
compatible pair has such an intersection.

\begin{theorem}[The relative-criterion comparison]\label{thm:criteria}
For every prime carrier atelier, with these exact conventions,
\[
\begin{gathered}
 \mathrm C\Rightarrow\mathrm D\Rightarrow\mathrm D_0,
 \qquad \mathrm D\Rightarrow\mathrm L3,
 \qquad \mathrm L2\Rightarrow\mathrm D_0,\\
 \mathrm L2+\mathrm L3\Rightarrow\mathrm D,
 \qquad \mathrm L3+\mathrm{CL}\Rightarrow\mathrm C,\\
 \mathrm D_0+\mathrm L1+\mathrm{CL}\Rightarrow\mathrm L2,
 \qquad \mathrm D_0+\mathrm L1+\mathrm{CL}_{\rm p}\Rightarrow\mathrm L2,\\
 \mathrm L2+\mathrm{CL}\Longleftrightarrow\mathrm L2+\mathrm{CD2},
 \qquad
 \mathrm L2+\mathrm{CL}_{\rm p}\Longleftrightarrow
                  \mathrm L2+\mathrm{CD2}_{\rm p}.
\end{gathered}
\]
Moreover $\mathrm C$ implies the two-parent matching condition:
for compatible $F,G$, refinements with $P|K=Q|K$ are compatible.
That matching condition is equivalent to unique gluing of
pairwise-matching refinements over every finite incidence cover.

The same gluing conclusion holds over arbitrary set-indexed covers
if the atlas admits unions of directed incidence-increasing
families $P_i\ll F$ whose restrictions to each principal stratum
of $F$ are eventually constant. Full canonical clique atlases
have the required closure. Finite gluing alone does not imply
this closure, even under all four lieu and absolute conditions
$\mathrm L1,\mathrm L2,\mathrm L3,\mathrm{CL}$.
\end{theorem}

\paragraph{Restricted intersection classes.}
Let $\Sigma$ be a class of posets closed under isomorphism and
containing every occurring principal deployment. The restricted
criterion is the following biconditional for elementary $x,y$,
with $K=x\cap y$:
\[
 x\bowtie y\quad\Longleftrightarrow\quad
 \left\{\begin{array}{l}
 (\forall p\in\ell(x)\setminus\ell(K),q\in\ell(y)\setminus\ell(K))
       \ p\parallel q,\\
 K\in\Sigma\quad\text{as an incidence poset}.
 \end{array}\right.
\]
\begin{theorem}[The exact intersection-class obstruction]\label{thm:sigma}
Assume the prime core, $\mathrm L3$, and this restricted criterion.
For compatible elementary parents $X,Y$, elementary refinements
$x\ll X,y\ll Y$, and compatible restrictions to $X\cap Y$,
\[
 x\bowtie y\quad\Longleftrightarrow\quad x\cap y\in\Sigma.
\]
Consequently relative compatibility $\mathrm C$ holds exactly
when every such common deployment belongs to $\Sigma$.
It holds for unrestricted $\Sigma$. If $\Sigma$ consists of the
empty poset and posets with a greatest element, it also holds
provided every elementary refinement restricts to an empty or
elementary figure on each elementary incidence subfigure of its parent.
Without that additional restriction hypothesis, a seven-stratum,
twenty-one-figure prime atelier satisfies $\mathrm L1,\mathrm L2,
\mathrm L3$, the restricted criterion, $\mathrm D$ and matching-cover
gluing, but fails $\mathrm C$.
\end{theorem}


\paragraph{Why these criteria work.}
Orthogonality is controlled by polar neighbourhoods. Splitting a poset point into a short interval preserves every such neighbourhood after closing, which brings Jen\v ca's quotient theorem into the problem. In the figure theory the missing property is primality: irreducibility by itself does not make a join into a union of deployed strata. For maquettes, a signed diamond expresses $\partial^2=0$, while a spherical boundary is the separate topological condition making the next cell a ball.

\section{Polar criteria}
We establish the general tests before specializing to intervals. Symmetry gives $A\subset N^2(A)$ and $N^3(A)=N(A)$. Thus $N^2$ is a closure operator and $N$ is an order-reversing involution on closed sets. Irreflexivity gives $A\cap N(A)=\varnothing$. These identities prove the asserted complete ortholattice structure, including the empty orthoset.

A maximal pairwise orthogonal subset $B$ of a closed set $A$ is called a basis. Dacey's criterion says that the lattice is orthomodular exactly when every basis of every closed set generates it. Here is the proof in the form used below. If the lattice is orthomodular, put $C=\cl B\le A$. Maximality gives $A\cap N(B)=\varnothing$, so the orthomodular law yields $A=C$. Conversely, extend a basis of a closed $C\le A$ to a basis of $A$, using Zorn's lemma. Every added element lies in $A\cap C^\perp$. If bases generate, it follows that $A=C\vee(A\wedge C^\perp)$. This proves the criterion.

For the Boolean test observe that $N(x)\cup N(y)\subset N(z)$ is equivalent to $z\in\cl\{x\}\cap\cl\{y\}$. In a Boolean algebra nonorthogonality implies nonzero intersection, so the condition is necessary. Conversely the condition implies that two closed sets with zero intersection are orthogonal: otherwise choose nonorthogonal $x,y$ in them and obtain a point in their intersection.

In an ortholattice this property implies distributivity. For $a\le b$, set $c=a\vee(b\wedge a^\perp)$. Then $b\wedge c^\perp=0$, so $b\le c$ and the orthomodular law holds. For arbitrary $a,b,c$, put $d=(a\wedge c)\vee(b\wedge c)$ and $e=(a\vee b)\wedge c$. The element $t=e\wedge d^\perp$ has zero meet with $a$ and with $b$, hence is below both $a^\perp,b^\perp$. It is also below $a\vee b$, so it is zero. Orthomodularity gives $e=d$, which is distributivity.

For the atom test, $y\in\cl\{x\}$ is equivalent to $N(x)\subset N(y)$. The closure of every such singleton is a nonzero closed subset of $\cl\{x\}$. Atomicity therefore forces its equality with $\cl\{x\}$, equivalently equality of polars. Conversely any nonzero closed subset contains a point $y$ and hence its singleton closure; the specified equality makes that subset equal to $\cl\{x\}$. This proves Theorem~\ref{thm:tests}.

The stronger equality $\cl\{x\}=\{x\}$ holds exactly when
$N(x)\subseteq N(y)$ implies $y=x$. Equality of the two polars
alone permits different points to generate the same atom.
For example, two points with no orthogonality have the two-element
Boolean support lattice and generate its same nonzero atom.
Under Proposition~\ref{prop:lieu-supports}, atomicity of a lieu
support can be tested entirely on the induced orthoset of lieux,
because trace sends its singleton closure to the corresponding
singleton closure there.

\section{Splitting points and removing boundedness}
For a poset $P$, let $Q^+(P)$ be its proper intervals, with
\[
 [u<v]\perp[r<s]\iff v\le r\text{ or }s\le u.
\]
Form $S(L)=L\times\{-,+\}$ with $x^-<x^+$ and with every copy of $x$ below every copy of $y$ when $x<y$. Embed $M(L)$ by
\[
 P_x\mapsto[x^-<x^+],\qquad I_{ab}\mapsto[a^+<b^-].
\]
Substitution shows that this is an induced orthogonality embedding. Call its image $U$.

Every interval of $S(L)$ is a concatenation $D(q)$ of at most three intervals of $U$: a possible first vertical interval, the middle interval between base points, and a possible last vertical interval. One has
\begin{equation}\label{eq:pieces}
 q\perp r\quad\Longleftrightarrow\quad(\forall d\in D(q))\ d\perp r.
\end{equation}
Indeed each piece must lie to one side of $r=[v<w]$. Switching sides at a common endpoint would force $w\le v$, so all pieces lie on the same side. For $B\subset U$ put $E(B)=\{q:D(q)\subset B\}$. Equation~\eqref{eq:pieces} gives $E(B)^\perp=E(B^\perp)$, with the ambient polar understood. Thus $E$ carries closed sets to closed sets and commutes with complement. Conversely a closed set contains an interval exactly when it contains all its pieces. Consequently trace on $U$ and $E$ are inverse complete ortholattice isomorphisms:
\[
 \Sigma(M(L))\simeq\Sigma(Q^+(S(L))).
\]
This comparison has no finiteness or boundedness hypothesis.

Suppose $P$ is a finite conditional lattice. Adjoin strict bounds $\widehat0,\widehat1$ to make a finite lattice $\widehat P$. For a closed $X\subset Q^+(P)$, close it in $Q^+(\widehat P)$ to obtain $\widehat X$. Then $\widehat X\cap Q^+(P)=X$: membership in the larger double polar implies membership in the smaller double polar, and the reverse containment is extensivity.

Jen\v ca's bounded theorem says that $\widehat X$ is a finite union of full quotient intervals arranged consecutively along a chain. Their traces are mutually orthogonal blocks, each a bounded interval block, a lower-ray block, an upper-ray block, or all of $Q^+(P)$. In a maximal orthogonal subset of a nonempty block, consecutive intervals have no gap: a gap would provide an extra orthogonal interval. They therefore concatenate. In a bounded block maximality forces the two extreme endpoints to be the endpoints of the block; its polar is exactly the polar of the whole block by \eqref{eq:pieces}. In a lower-ray block the first endpoint is minimal in all of $P$ and the last is the ray bound, so its polar consists precisely of intervals above that bound. The upper-ray case is dual. In the whole-space case maximality means its polar is empty. Hence a maximal orthogonal subset of $X$ has the same polar as $X$. Dacey's criterion proves orthomodularity.

The split poset $S(L)$ is a conditional lattice exactly when $L$ is: for incomparable base points, a meet lifts to $(x\wedge y)^+$ and a join to $(x\vee y)^-$; comparable copies have their usual minimum and maximum. The comparison proves sufficiency in Theorem~\ref{thm:interval}.

For necessity suppose $x,y$ have upper bounds but no join. Let $U$ be their common upper bounds and let $I$ be the common lower bounds of $U$. Then $I^\uparrow=U$, $U^\downarrow=I$, and $I\cap U=\varnothing$. Choose incomparable maximal $a,b\in I$ and a minimal element $m$ of $L$ below $a$. In $Q^+(S(L))$ put $X=Q^+(S(I))$ and $Y=Q^+(S(U))$. They are mutual polars. Indeed $Y\subseteq X^\perp$, since every member of $I$ lies strictly below every member of $U$. Conversely start with $q\in X^\perp$. It is orthogonal to every vertical interval over $I$. If it lies below one of them, both endpoint bases belong to the downset $I$, so $q\in X$; this contradicts $q\in X^\perp$ and irreflexivity. Therefore it lies above every vertical interval over $I$, placing both endpoint bases in $I^\uparrow=U$. Thus $X^\perp=Y$. The order-dual argument gives $Y^\perp=X$.

The singleton $B=\{[m^-<a^+]\}$ is maximal orthogonal in $X$: no interval lies below $m^-$ or above $a^+$ within $S(I)$. For $u\in U$, however, $[a^+<u^-]$ is orthogonal to $B$ and is not orthogonal to $[b^-<b^+]$. Thus $B$ does not generate $X$, violating Dacey's criterion. A missing bounded meet is handled by reversing the order. This proves the finite classification.

For Booleanity in arbitrary $L$, we first show that every nonisolated $x$ has singleton closure $\cl\{P_x\}=\{P_x\}$. A point $y\ne x$ comparable to $x$ excludes itself: $P_y\in N(P_x)\setminus N(P_y)$. For $y$ incomparable to $x$, choose $z$ comparable to $x$. The interval between $x,z$ is orthogonal to $P_x$, but cannot be orthogonal to $P_y$, as either separation inequality would make $x,y$ comparable. Thus no other point is in the closure.

If an interval $I_{ab}$ were in that closure, its orthogonality to the same incident interval would force $x\le a$ or $b\le x$. In the first case $I_{xb}$ is orthogonal to $P_x$ but not to $I_{ab}$; in the second case use $I_{ax}$. Both contradict the double-polar condition. This proves the singleton assertion. Distinct singleton atoms in a Boolean ortholattice are orthogonal, so all nonisolated points must be comparable.

Conversely let the nonisolated core $C$ be a nonempty chain. Form a set $E_C$ with a point atom for every $x\in C$ and a gap atom for every cover pair $a<b$ in $C$. Give $P_x$ its singleton point support when $x\in C$. Give $I_{ab}$ all point atoms strictly between $a,b$, together with all gap atoms whose closed endpoints lie in $[a,b]$. Give every isolated point symbol the full support $E_C$. Every such support is nonempty: an interval either has an interior point or is itself a cover pair.

Two symbols are orthogonal exactly when these supports are disjoint. For two intervals, if their endpoints fail both separation inequalities, put $r=\max(a,c)$ and $s=\min(b,d)$; then $r<s$. Either there is a point strictly between $r,s$, yielding a common point atom, or $(r,s)$ is a cover pair, yielding a common gap atom. The reverse implication and the point cases follow directly from the definitions. Every singleton support is represented, including a gap atom by its cover interval. Therefore the polar of any family has union of supports equal to the complement of that family's union of supports. Double polars are exactly the families of all symbols supported in a given subset of $E_C$. All subsets occur, so $\Sigma(M(L))\simeq\mathcal P(E_C)$.

If $C$ is empty and $L$ is nonempty, the orthoset has no edges and its support lattice has two elements. If $L$ is empty it has one element. We include the trivial Boolean algebra, so both cases satisfy the statement. This completes Theorem~\ref{thm:interval}.

\begin{corollary}[Atomic point supports]\label{cor:interval-atoms}
Every point symbol of $M(L)$ has atomic support if and only if
$L$ is an antichain or has no isolated points. The assertion is
vacuous for the empty poset.
\end{corollary}
\begin{proof}
The proof above gives a singleton closed support for each nonisolated
point. An isolated point has empty polar and hence full support.
If there is any comparable pair, a nonisolated point gives a proper
nonzero singleton support, so the full support is not an atom.
If there are no comparable pairs, the nonempty orthoset has its
two-element support lattice and every point generates its unique
nonzero atom. These are all cases.
\end{proof}

\subsection{The strong interval structure}
We prove Theorem~\ref{thm:interval-system} by checking the
factorization presentation. Incidence is a partial order because
its only proper relations run from an endpoint point to an
interval. Internal refinement is a partial order by endpoint
containment: the nonidentity transitivity cases are
$P_x\rho I_{cd}\rho I_{ab}$ and
$I_{ef}\rho I_{cd}\rho I_{ab}$, and their endpoint inequalities
give the required composite. Mutual interval containment forces
equal endpoints; a proper point-to-interval relation cannot reverse.

For a target point, an internal factor is that point itself.
For $I_{ab}$, the possible factors are $P_a,P_b,I_{ab}$.
A point source $P_x$ uses $P_a$ when $x=a$, $P_b$ when $x=b$,
and $I_{ab}$ when $a<x<b$; these cases are disjoint.
An interval source has only the candidate $I_{ab}$, which works
exactly under the stated containment inequalities. Factorization
is therefore unique when it exists, and gives the displayed
overall refinement.

For exchange, the only proper source incidence is an endpoint
of $I_{cd}$. If $I_{cd}\rho I_{ab}$, its endpoint is either
one of $a,b$, whose point is its factor, or strictly between them,
whose factor is $I_{ab}$. This proves exchange in every case.

The positional relation is a strict partial order. Transitivity
uses
$\mu(U)\leq\lambda(V)\leq\mu(V)\leq\lambda(W)$.
A positional cycle would force all its endpoints equal, and
hence all its symbols to be the same point, contradicting
strictness. Its comparability relation is exactly the previously
defined orthogonality.

If $X'\rho X<_{\mathrm p}Y$, endpoint containment gives
$\mu(X')\leq\lambda(Y)$. Equality of the symbols $X'=Y$
would force $Y$ to be the point at the right endpoint of $X$.
Such a point cannot internally refine a distinct interval; if
$X$ is a point it would already equal $Y$. Thus $X'\ne Y$
and $X'<_{\mathrm p}Y$. The left-endpoint argument handles
$Y<_{\mathrm p}X$, and symmetry gives heredity in both variables.
The principal deployment of an interval is
$P_a<_{\mathrm p}I_{ab}<_{\mathrm p}P_b$, so all its distinct
pairs are disjoint; a point deployment is a singleton.
All strong factorization axioms are proved, and
Theorem~\ref{thm:strong-relations} gives the carrier presentation.

Primitive cliques are precisely positional chains. Incidence
closure is precisely endpoint closure. Adjoining an endpoint
preserves such a chain: every symbol before $I_{ab}$ ends at or
before $a$, and every symbol after it starts at or after $b$.
Finite chains acquire at most two points per interval. This proves
the full and finite atlas assertions.

In any strong system a common internal refinement precludes
disjunction, by heredity and irreflexivity. If $L$ is a chain,
two nondisjoint points coincide. A point nondisjoint from $I_{ab}$
lies strictly between $a,b$ and is itself a common internal
refinement. Two nondisjoint intervals satisfy
\[
 \max(a,c)<\min(b,d),
\]
and the interval with these endpoints internally refines both.
This last witness is an interval symbol; no intervening point
or density assumption is needed. Thus every chain is faithful.
If $x,y$ are incomparable, $P_x,P_y$ are nondisjoint but have
no common internal refinement, so faithfulness fails.

Only $P_x$ refines the point $P_x$, whereas every interval has
a proper refining endpoint. Hence the lieux are exactly the
points, even for an interval with no internal point. The displayed
open-lieu formulas and their nonemptiness criterion follow directly.
In particular Corollary~\ref{cor:interval-atoms} is exactly the
all-lieux atomicity criterion. Every nonempty connected comparability
poset satisfies it: either it is a singleton or it has no isolated
points.

\section{Boolean calculus within a figure}
We give a criterion that distinguishes the calculus of one figure from global Booleanity. Let $(U_i)_{i\in I}$ be nonempty mutually orthogonal subsets of an orthoset $E$, and put $S_A=\cl(\bigcup_{i\in A}U_i)$ and $S=S_I$.

\begin{proposition}\label{prop:figure}
The map $A\mapsto S_A$ is always injective and preserves arbitrary joins. It has the relative-complement identities
\[
 S_A\cap N\Bigl(\bigcup_{i\in B}U_i\Bigr)=S_{A\setminus B}
 \quad(B\subseteq A\subseteq I)
\]
if and only if
\[
 z\in S,\quad (\forall i)\,[z\perp U_i\text{ or }w\perp U_i]
 \quad\Longrightarrow\quad z\perp w.
\]
When this holds it also preserves arbitrary nonempty meets; empty meets are preserved into $[0,S]$.
\end{proposition}
\begin{proof}
Joins follow from the closure identities. If $i\in A\setminus B$, a point of $U_i$ belongs to $S_A$ and is orthogonal to $S_B$, proving injectivity. For the displayed implication, let $z\in S_A\cap N(U_B)$ and $w\in N(U_{A\setminus B})$, with unions abbreviated. For indices outside $A$, $z$ is orthogonal to $U_i$ by mutual orthogonality and closure. The implication therefore gives $z\perp w$, proving the difficult containment in the complement identity. The other containment follows from orthogonality of different blocks.

Conversely put $B=\{i:z\perp U_i\}$. The complement identity for $A=I$ gives $z\in S_{I\setminus B}$, while the hypothesis puts $w\in N(U_{I\setminus B})$. Hence $z\perp w$. Finally $S_A=S\cap N(U_{I\setminus A})$; intersecting these formulas and using the polar of a union gives the arbitrary-meet assertion.
\end{proof}
An incidence-closed subfigure has the intended whole support precisely when its deployed set and the union of its open strata have the same polar. This is an additional coverage condition if the incidence axioms do not already imply it.

\section{Admissible figures and complete support maps}
\label{sec:figure-comparisons}

\subsection{Normality and its finite interval consequence}
We first prove Theorem~\ref{thm:flags}. Apply
Proposition~\ref{prop:figure} to the singleton blocks of a clique.
It says that admissibility of $F$ is equivalent to
\begin{equation}\label{eq:clique-test}
 z\in\cl F,\quad F\subseteq N(z)\cup N(w)
 \quad\Longrightarrow\quad z\perp w.
\end{equation}
A maximal clique has empty polar and hence closure $E$. Thus
admissibility of every clique gives the stated maximal-clique
condition. Conversely extend $F$ to a maximal clique $D$. Every
$d\in D\setminus F$ lies in $N(F)$ and is orthogonal to any
$z\in\cl F$. Therefore the premise in \eqref{eq:clique-test}
implies $D\subseteq N(z)\cup N(w)$. The case $z=w$ is impossible:
it would put $z$ in both $N(F)$ and $\cl F$. The maximal-clique
condition gives $z\perp w$, as required. This proves the equivalence,
including infinite cliques when maximal extensions are chosen.

For the finite interval assertion, give a point symbol equal left
and right endpoints and give $I_{ab}$ endpoints $a,b$. Define the
positional order by
\[
 Z<_{\mathrm p}W\quad\Longleftrightarrow\quad
 Z\ne W\ \hbox{ and the right endpoint of $Z$ is at most the
 left endpoint of $W$}.
\]
It is a strict partial order, and its comparable pairs are exactly
the orthogonal pairs. A finite orthogonal subset is therefore a
finite positional chain. Its endpoints form a chain
$t_0<\cdots<t_n$ in $L$. Insert every $P_{t_i}$ and each
$I_{t_it_{i+1}}$ to obtain an alternating chain $H$. It contains
the original subset: an original interval cannot contain another
endpoint strictly between its endpoints, since the corresponding
symbol would not be orthogonal to it.

For any symbol $Z$, let $T_Z$ consist of the members of $H$ not
orthogonal to $Z$. If nonempty, this is a contiguous block.
Indeed, a member between two incomparable members can be neither
below nor above $Z$, by transitivity. Every member before the block
is below $Z$, and every member after it is above $Z$.
Suppose two nonempty blocks $T_Z,T_W$ are disjoint, with $T_Z$
first. At a consecutive cut $h_i,h_{i+1}$ between them,
$Z<_{\mathrm p}h_{i+1}$ and $h_i<_{\mathrm p}W$.
The consecutive symbols are $P_t,I_{tu}$ or $I_{st},P_t$.
Thus the right endpoint of $Z$ is at most $t$, and the left
endpoint of $W$ is at least $t$. It follows that
$Z<_{\mathrm p}W$. The only possible failure of strictness would
be the same point symbol twice, which the disjoint nonempty
blocks exclude.

Now take $Z\in\cl H$ and $H\subseteq N(Z)\cup N(W)$.
The first trace is nonempty, and the two traces are disjoint.
If $T_W$ is empty, $W\in N(H)$ gives $Z\perp W$ directly.
Otherwise the preceding cut argument gives the same conclusion.
Equation~\eqref{eq:clique-test} proves admissibility of $H$.
For a subchain $F\subseteq H$ and $A\subseteq F$,
$\cl F\subseteq N(H\setminus F)$, whence
\[
 \cl F\cap N(A)
 \subseteq\cl H\cap N((H\setminus F)\cup A)
 =\cl(F\setminus A).
\]
The reverse inclusion is immediate. Empty subsets are included.
This proves the finite assertion in Theorem~\ref{thm:flags}.

\begin{counterexample}[Finite figure complements do not force orthomodularity]
\label{ex:eight-symbols}
Let $L=\{a,b,c,d\}$ with exactly
$a<c,a<d,b<c,b<d$. The eight-symbol orthoset $M(L)$ has twenty
supports and is normal, but its support lattice is not orthomodular.
\end{counterexample}
\begin{proof}
There are four point edges and, for each interval symbol, the two
edges joining it to its endpoints. The supports are the empty set,
the eight singletons, the six pairs of point symbols, the four
point neighborhoods, and the whole set.
Here $N(I_{xy})=\{P_x,P_y\}$, while intersecting the neighborhoods
of an upper and a lower endpoint gives their interval singleton.
Intersecting the two upper point neighborhoods gives
$\{P_a,P_b\}$, and dually for the two lower neighborhoods.
Further intersections give the point singletons. The displayed
family is closed under intersections and contains all point and
interval neighborhoods, so it is precisely the family of all
polars, establishing the count.

The closed sets $A=\{P_a\}\subset B=\{P_a,P_b\}$ satisfy
$B\cap N(A)=\varnothing$, contradicting the orthomodular law.
The maximal cliques are the four endpoint--interval triangles.
If such a triangle lies in $N(u)\cup N(v)$, its interval forces
one of $u,v$ to be an endpoint. Covering that endpoint forces the
other vertex to be its neighbor. Thus the maximal-clique condition
holds. The finite-flag result proves the same figure-complement
conclusion without the enumeration.
\end{proof}

\begin{counterexample}[An infinite flag across a cut]\label{ex:infinite-flag}
Let $L=\Q\amalg\{a,b\}$, with the usual rational order,
$a<q$ exactly for $q\ge0$, and $q<b$ exactly for $q<0$.
There are no further comparisons involving $a,b$.
The orthogonal family $F=\{P_q:q\in\Q\}$ is not admissible.
\end{counterexample}
\begin{proof}
Its polar is empty. A rational point fails at itself; $P_a$ misses
negative rationals and $P_b$ misses nonnegative rationals. An
interval with rational endpoints misses an interior rational.
An interval starting at $a$ misses negative rationals, and one
ending at $b$ misses nonnegative rationals. These exhaust the
possible symbols. Thus $\cl F=M(L)$.
Put $A=\{P_q:q\ge0\}$ and $B=F\setminus A$.
Then $P_a\in N(A)$ and $P_b\in N(B)$, but $P_a\not\perp P_b$.
Consequently $P_a\notin\cl B$ although
$P_a\in\cl F\cap N(A)$.
\end{proof}

\subsection{Which families detect the orthomodular law?}
\begin{proposition}\label{prop:orthogonal-families}
In a complete orthomodular lattice, every mutually orthogonal
family $(e_i)_{i\in I}$ of nonzero elements gives an injective
complete Boolean embedding into $[0,\bigvee_i e_i]$ by
$A\mapsto\bigvee_{i\in A}e_i$.
Conversely, the relative-complement identity for every two-element
orthogonal family of nonzero elements implies orthomodularity.
\end{proposition}
\begin{proof}
For orthogonal $u,v$, the dual orthomodular identity gives
$(u\vee v)\wedge u^\perp=v$. Apply this with
$u=\bigvee_{i\in A}e_i$ and $v=\bigvee_{i\notin A}e_i$.
These joins are orthogonal, even for an infinite family, by
completeness and reversal of order under complement.
The relative complement follows. Joins are preserved directly;
intersections follow by complementing unions inside the displayed
interval. An index in $A\setminus B$ contributes a nonzero element
below the first join and orthogonal to the second, so the map is
injective.

Conversely let $a\le b$. Use the two orthogonal elements
$u=b^\perp$ and $v=a$, omitting a zero term. The assumed identity
gives $(b^\perp\vee a)\wedge b=a$. Complementing gives the
orthomodular law for $b^\perp\le a^\perp$, hence for every
comparable pair. Counterexample~\ref{ex:eight-symbols} shows why
testing only the families supplied by figures is weaker.
\end{proof}

\begin{proposition}[Faithful refinement and regular opens]\label{prop:regular-open}
Let $\rho$ be a preorder on $E$ and define $x\perp y$ to mean that
there is no $z$ with $z\rho x$ and $z\rho y$. Then $\Sigma(E)$
is the complete Boolean algebra of regular opens of the lower
Alexandrov topology of $\rho$.
\end{proposition}
\begin{proof}
Write $O_A=\bigcup_{a\in A}\mathord\downarrow_\rho a$.
A point is orthogonal to $A$ exactly when its principal lower set
misses $O_A$. Therefore
\[
 N(A)=\operatorname{int}(E\setminus O_A),\qquad
 \cl A=\operatorname{int}\overline{O_A}.
\]
The latter is regular open, and every regular open $U$ is obtained
by taking $A=U$. Two disjoint regular opens are orthogonal because
a common lower refinement of points in them would lie in their
intersection. The Boolean argument in Section~2 applies.
Arbitrary intersections of regular opens are open in an Alexandrov
topology and remain regular by monotonicity of
$\operatorname{int}\overline{(-)}$. These are the complete meets;
joins are regular-open closures of unions and complement is
interior of the set-theoretic complement.
\end{proof}

\begin{proposition}[A sufficient graph class with a strict converse]
\label{prop:cographs}
Every finite cograph has an orthomodular support lattice. A
five-vertex comparability graph can have a Boolean support lattice
without being a cograph.
\end{proposition}
\begin{proof}
Define cographs recursively from a one-vertex graph by disjoint
union and graph join. For a graph join, taking polars acts
coordinatewise on the vertex parts, so its support lattice is the
product of the support lattices of those parts. For a disjoint
union, every proper nonzero closed set lies in one part and is
proper and closed there. Complements agree with the complement
in that part. Thus the support lattice is the horizontal sum,
identifying all zeros and all tops; a two-element factor adds no
proper element. Products and horizontal sums preserve
orthomodularity. For the latter, every comparable proper pair
lies in a single factor. Induction proves the first assertion.

For strictness use five points represented by
$\{1\},\{2\},\{3\},\{1,2\},\{2,3\}$, orthogonal when
disjoint. The three singleton points detect the union represented
by any subset of points. Its double polar consists exactly of
the five points whose representing subsets lie in that union.
All eight unions occur, so the support lattice is the Boolean
algebra on three atoms. Calling the points $a,b,c,x,y$ in the
displayed order, the graph has edges $ab,ac,bc,cx,ay$.
It is the comparability graph of $a<b<c$, $x<c$, $a<y$.
The path $x-c-a-y$ is induced, which a cograph cannot contain:
an induced four-vertex path is neither a nontrivial disjoint union nor a
nontrivial join, and the recursive construction is inherited by induced
subgraphs. This proves strictness.
\end{proof}

\subsection{Restriction to a set of points}
The following criterion identifies exactly when a smaller orthoset
determines every support. Let $L\subseteq M$ have the orthogonality
induced from an orthoset $M$. Suppose each $x\in M$ has a nonempty
set $U_x\subseteq L$, with $U_p=\{p\}$ for $p\in L$, such that
\begin{equation}\label{eq:lieu-detection}
 x\perp y\quad\Longleftrightarrow\quad U_x\perp U_y.
\end{equation}
Orthogonality on the right means every pair is orthogonal.

\begin{proposition}[Exact support restriction]\label{prop:lieu-supports}
Under these hypotheses, the maps
\[
 S\longmapsto S\cap L,\qquad
 T\longmapsto\widehat T=\{x\in M:U_x\subseteq T\}
\]
are inverse complete ortholattice isomorphisms
$\Sigma(M)\simeq\Sigma(L)$.
Consequently the orthomodular and Boolean criteria are exactly
the corresponding criteria on $L$. Every orthoset is a possible
point orthoset in this statement.
\end{proposition}
\begin{proof}
Put $U(A)=\bigcup_{x\in A}U_x$. Detection gives
\[
 N_M(A)=\widehat{N_L(U(A))},\qquad
 \widehat T\cap L=T.
\]
If $T$ is closed in $L$, then
$N_M(\widehat T)=\widehat{N_L(T)}$: one implication tests the
points $p\in T\subseteq\widehat T$, and the other uses detection
and $U_y\subseteq T$ for $y\in\widehat T$. Applying this twice
proves that $\widehat T$ is closed in $M$. Conversely a closed
support $S=N_M(A)$ is the hat of the closed polar $N_L(U(A))$.
The trace formula proves that these maps are inverse; the polar
formula proves that they preserve orthocomplements. An order
isomorphism preserves all joins and meets. Taking $M=L$ and
$U_x=\{x\}$ realizes any given orthoset, including the empty one.
\end{proof}

In a refinement presentation the points called \emph{lieux} are
the minimal refinement objects, and $U_x$ is the set of lieux
internally refining $x$. Minimality gives $U_p=\{p\}$.
Thus nonempty open lieu sets and the explicit detection
condition~\eqref{eq:lieu-detection} are the precise inputs of this
support comparison. No additional restriction on the point graph
follows from the general formalism: incidence and refinement can
both be identity, with the given orthogonality, so every object
is a lieu and each principal deployment is a singleton.

\paragraph{Faithfulness and lieu atoms.}
For a strong system with nonempty open-lieu sets and local detection,
the following conditions are equivalent: faithfulness; pairwise
disjunction of distinct lieux; and Booleanity of its support lattice
with distinct lieux generating distinct atoms.
Indeed, distinct lieux have no common internal refinement, so
faithfulness forces their disjunction. Conversely, if the lieu
graph is complete, local detection says that two strata are
nondisjoint exactly when their open-lieu sets intersect. A common
lieu is a common internal refinement; any common internal
refinement contains a lieu by nonemptiness and internal
transitivity. This is faithfulness.
The induced support lattice on a complete lieu graph is
$\mathcal P(L)$, and Proposition~\ref{prop:lieu-supports} gives
the corresponding statement for all strata. Conversely distinct
atoms of a Boolean ortholattice are orthogonal, so distinct
lieux must be disjoint.

If the lieu graph additionally has singleton closed supports,
Booleanity alone forces distinct lieux to be disjoint: two
different singleton atoms in a Boolean algebra are orthogonal.
Without that separation condition, two nondisjoint lieux may
generate the same Boolean atom, as the two-point example in
Section~2 shows. For a connected comparability poset $L$, the
nonisolated-core criterion in Theorem~\ref{thm:interval} reduces
to totality of $L$, with the singleton case included.

The separation issue also occurs with genuine interval strata.
Take $L=\mathbb Q\amalg\{c\}$, with $c$ incomparable to every
rational. Its nonisolated core is the chain $\mathbb Q$, so its
support lattice is Boolean. Every rational interval has a nonempty
open-lieu set by density, and all point open-lieu sets are
singletons. Local detection holds: two rational interval interiors
overlap exactly when they have a common rational point, which
witnesses failure of disjunction; the point--interval case is the
same membership test. Distinct rational points are disjoint.
The point $c$ is disjoint from no rational point, and every open
lieu set being tested is nonempty, so detection also holds for
pairs involving $c$. Nevertheless $P_c$ is nondisjoint from every
other point and has full support because its polar is empty.
Thus nonempty open-lieu sets and Booleanity do not replace the
stated separation condition.

\subsection{Objects and correspondences}
We prove Theorem~\ref{thm:correspondence} in full. On
$E_K=K\setminus\{0\}$ the relation $x\le Cy$ is symmetric and
irreflexive. For $A\subseteq E_K$ its polar is the nonzero
principal ideal below $C(\bigvee A)$, and its double polar is the
nonzero principal ideal below $\bigvee A$. Every principal ideal
occurs, and inclusion and complement give the asserted complete
ortholattice isomorphism. The empty orthoset represents the
one-element lattice.

The zero convention has a precise variant. Suppose a symmetric
relation permits self-orthogonal elements only when they are
orthogonal to every element, and let $R$ be their set. Every polar
contains $R$. Removing $R$ gives an ordinary orthoset, and
$S\mapsto S\setminus R$, with inverse adjoining $R$, identifies
the two complete support lattices. Thus retaining zero does not
silently change the irreflexive convention.

If $f$ preserves all joins, then
$\cl A=\bigvee_{a\in A}\cl\{a\}$ proves the correspondence
condition for $\alpha(x)=f(\cl\{x\})$. Conversely that condition
gives
\[
 \bigvee_{x\in\cl A}\alpha(x)=\bigvee_{a\in A}\alpha(a).
\]
It follows immediately that $f(S)=\bigvee_{x\in S}\alpha(x)$
preserves all joins. Taking $A=\{x\}$ and using
$x\in\cl\{x\}$ shows that these assignments are inverse.
Substitution gives the identity and composition formulas.
The object representation just proved makes the resulting
full and faithful functor essentially surjective.

Every complete join map has right adjoint
$g(T)=\bigvee\{S:f(S)\le T\}$. Hence
\[
 f(S)\le C_FT\quad\Longleftrightarrow\quad
 S\le g(C_FT)\quad\Longleftrightarrow\quad
 S\perp C_Eg(C_FT).
\]
The conjugate $C_EgC_F$ preserves joins since $g$ preserves meets.
The displayed characterization proves uniqueness, reverses
composition, and gives $(f^\dagger)^\dagger=f$.
Equivalently the correspondence condition says that each column
$\{x:\alpha(x)\perp y\}$ is closed: apply the condition to
that column for one direction, and test all elements of the
polar of $\bigvee_{a\in A}\alpha(a)$ for the other.

Finally, for a classical orthoset, where distinct points are
orthogonal, supports are all subsets. Complete join maps between
such powersets are arbitrary relations. A relation sending one
point to two is not a function, and a function identifying two
points need not preserve orthogonality. Both give the specified
correspondences. This proves the stated distinction of categories.

\subsection{Strong relational systems and clique atlases}
We prove Theorem~\ref{thm:strong-relations}, including the
partial-order assertions for infinite figures.

Start with the carrier presentation. Internal refinement is
reflexive and antisymmetric because refinement is a partial order.
If $c=c(x,y)$, then $x\rho c$: any smaller candidate below $c$
would also contradict its leastness below $y$.
If $x\rho z\leq y$, then $c(x,y)\leq z$, and the carrier
condition below $z$ forces $c(x,y)=z$. This proves uniqueness
of factorization. Applying the carrier construction to
$x'\leq x\rho y$ proves exchange.

It remains to justify internal transitivity. If $x\rho y\rho z$,
write $x\rho c\leq z$ for its carrier factorization.
If $c\ne z$, principal disjointness gives $c\delta z$.
Heredity along $y\rho z$ and $x\rho y$ gives $c\delta x$;
heredity along $x\rho c$ then gives $x\delta x$, a contradiction.
Thus $c=z$ and $x\rho z$.

Conversely start with the factorization presentation and define
$\ll$ as in the theorem. Reflexivity and containment of incidence
are immediate. For transitivity, from
$x\rho a\leq y\rho b\leq z$, exchange gives
$a\rho c\leq b$, and then $x\rho c\leq z$.
For antisymmetry suppose also $y\rho b\leq x$.
The same exchange gives $x\rho a\rho c\leq x$.
Uniqueness of factorization into $x$, compared with
$x\rho x\leq x$, gives $c=x$. Antisymmetry of $\rho$
gives $a=x$, hence $x\leq y$. Reversing the roles gives
$y\leq x$, so $x=y$.
If $x\rho c\leq y$ and $x\ll d\leq y$, factorization
through $d$ gives the same $c$, so $c\leq d$.
Thus $c$ is the least carrier and the recovered internal relation
is exactly the original one. The constructions are inverse.

Now let $\mathscr H$ be an admitted clique atlas as in the theorem.
Bounded unions are admitted by downward-closed subfigure closure
inside their bound. Every figure is the union of its principal
downsets. Those principals are completely join-prime by membership
of their maximum, and this union expression shows that they are
exactly the nonzero completely join-irreducible figures.
Two figures are compatible exactly when their union is a clique.
For principal downsets this says
\[
 x\bowtie y\quad\Longleftrightarrow\quad
 (\forall a\leq x,\ a\not\leq y)
 (\forall b\leq y,\ b\not\leq x)\ a\delta b.
\]
Pairs lying in one principal downset already satisfy the required
disjunction; these are precisely the remaining pairs.

A stratum $p$ refining a figure $F$ has a carrier: factor its
refinement to some $x\in F$ as $p\rho c\leq x$.
Two possible carriers $c,d\in F$, if distinct, would satisfy
$c\delta d$, and heredity would imply $p\delta p$.
They therefore agree and give the least refined stratum in $F$.
If $P,F$ are compatible and $P\ll F$, the carrier $c$ of
$p\in P$ lies in their common clique. If $p\ne c$,
$p\delta c$ and $p\rho c$ again give a self-disjunction.
Thus $p=c\in F$ and $P\subseteq F$.

Elementwise refinement is reflexive and transitive. Suppose
$P\ll F\ll P$. For $p\in P$, take its successive carriers
$q\in F,r\in P$. Then $p\rho q\rho r$.
The carrier of $p$ in $P$ is $p$, so internal transitivity and
uniqueness give $r=p$. Antisymmetry of $\rho$ yields $q=p$.
Hence $P\subseteq F$; symmetry gives equality. This proves
refinement antisymmetry without finiteness. The remaining
carrier-atelier clauses now follow directly from the definitions.

\begin{proposition}[The abstract restriction calculus]
\label{prop:abstract-carriers}
In any prime carrier atelier, carriers compose. If $P\ll F$
and $G\leq F$, the strata of $P$ whose $F$-carriers lie in
$\widetilde G$ form a figure $P|G$. It is the greatest common
lower bound of $P,G$ for refinement. Within one parent,
restriction commutes with incidence unions and intersections.
\end{proposition}
\begin{proof}
Prime generation identifies the incidence ideal below a figure
with all downward-closed subsets of its stratum deployment:
the join of such a subset has exactly that deployment by primality.
If $x\ll P\ll F$, put $y=c_P(x)$ and $z=c_F(y)$.
Then $x\rho y\rho z$, so uniqueness gives $z=c_F(x)$.
If $x'\leq x$, the carrier of $x$ is a candidate refined by
$x'$, whence $c_F(x')\leq c_F(x)$. Thus the specified inverse
image of $\widetilde G$ is downward closed and defines a figure.
It refines both $P$ and $G$.

If $Q\ll P,G$ and $q\in\widetilde Q$, let $p=c_P(q)$.
Composition gives $c_F(p)=c_F(q)$. Since $q\ll G\leq F$,
the latter carrier lies in $\widetilde G$.
Thus $p\in\widetilde{P|G}$ and $q\ll p$; the stratum
test gives $Q\ll P|G$. This proves the infimum property.
The union and intersection assertions are the corresponding
inverse-image identities for the carrier map.
\end{proof}

For the clique atlas, let $F,G$ be compatible parents,
$K=F\cap G$, and $P\ll F,Q\ll G$.
Compatibility of $P|K,Q|K$ implies compatibility of $P,Q$.
Indeed, for distinct $p\in P,q\in Q$, compare their parent
carriers $c,d$. If $c\ne d$, the common parent clique gives
$c\delta d$, hence $p\delta q$. If $c=d$, it belongs to
$K$, and the required disjunction is provided by the two
compatible restrictions.
If those restrictions are also disjoint as sets, then $P,Q$
are disjoint: a common stratum would have two parent carriers;
distinct ones force self-disjunction, while equal ones put it
in both restrictions. In particular disjoint parents have
disjoint refinements. Distinct compatible strata likewise have
pairwise disjoint open lieux, by heredity along their internal
refinements.

For the gluing statement of Theorem~\ref{thm:strong-relations},
the matching restrictions make every pair $P_i,P_j$ compatible
by the preceding argument. Their union $P$ is admitted for a
finite cover, and for any cover in the full atlas. It refines $F$.
If $p\in P_j$ has carrier in $F_i$, that carrier lies in
$F_i\cap F_j$; matching puts $p$ in $P_i$.
Therefore $P|F_i=P_i$. Conversely every stratum of any proposed
refinement of $F$ is carried by some covering $F_i$, so these
restrictions determine it. This proves uniqueness.
A directed incidence union is a clique downset because any two
of its strata lie together in one member, proving the final
full-atlas assertion.

\subsection{The representation by admissible set figures}
We prove Theorem~\ref{thm:set-atlas} in both directions.
Start with a prime carrier atelier and use its canonical sets
\[
 A_x=\{z:z\ll x\},\qquad \Phi_F=\{A_x:x\leq F\}.
\]
One has $A_x\subseteq A_y$ if and only if $x\ll y$;
thus the sets distinguish the strata. For a point $z$ in the
support of $\Phi_F$, its least containing member is
$A_{c_F(z)}$. The cell of $A_x$ is consequently
\[
 U_x=\{z:z\rho x\},
\]
independent of the parent figure and nonempty because $x\rho x$.
Hence $\Phi_F$ is admissible.
Inside one figure, inclusion $A_x\subseteq A_y$ is compatible
refinement and therefore means $x\leq y$. The principal
subfigure at $A_x$ is exactly $\Phi_x$.

Incidence is preserved and reflected by prime deployment.
If $F\ll G$ and $x\leq F$ has carrier $y$ in $G$,
internal transitivity gives $U_x\subseteq U_y$, proving
set-figure refinement. Conversely $U_x\subseteq U_y$
puts $x\in U_y$ and hence $x\rho y$; the stratum test
recovers the original refinement.

Prime generation turns bounded joins into unions, giving
condition~1. The principal identity gives condition~2.
Inclusion $A_x\subseteq A_y$ is refinement of $x$ to $y$,
so the preceding argument gives condition~3. Compatibility
detected on strata gives condition~4. Finally the global
family is all principal refinement ideals $A_x$. A point $x$
has least member $A_x$, and that member has global residual
$\{x\}$. These sets cover $M$ and separate points by
antisymmetry of refinement. This proves condition~5.

Conversely suppose the five set-atlas conditions hold.
Global admissibility gives each point $e$ a least member $A_e$.
It belongs to a global stratum $B$ precisely when $A_e\subseteq B$.
Two points with the same least member therefore have the same
memberships; separation makes them equal. Nonempty residuals
give a unique point $e_A$ for every occurring stratum $A$.
Thus points and occurring strata are in bijection.

Order $\mathscr H$ by inclusion. Bounded unions exist by
condition~1. By condition~2 every figure is the union of its
admitted principal subfigures $\widehat A$. Each is completely
join-prime: membership of its maximum $A$ in a union occurs
in one summand, which then contains all of $\widehat A$.
The union expression shows that these are exactly the nonzero
completely join-irreducibles. Condition~4 gives compatibility
detected on strata. A closed subfamily is admitted because it
is a bounded union of its principal subfigures.

Principal consistency makes the cell $C_A$ of $A$ independent
of its admitted parent. Set-figure refinement is reflexive and
transitive. If two figures refine each other, their nonempty
partition cells coincide. Equality of a cell of $A$ with one
of $B$ puts $e_A\in B$ and $e_B\in A$, giving both
$A\subseteq B$ and $B\subseteq A$. Thus $A=B$ and
the figures coincide. Refinement is a partial order. Inclusion
of figures preserves their cells and hence implies refinement.

If a cell of $A$ is contained in a cell of $B$, then
$e_A\in B$, so $A\subseteq B$. Condition~3 gives
$\widehat A\ll\widehat B$. Conversely, refinement of
$\widehat A$ to a principal subfigure of $\Psi$ maps its
top cell into a cell of $\Psi$. This proves exactly the
elementwise refinement axiom.

For $\widehat A\ll\Psi$, let $B$ be the least member of
$\Psi$ containing $e_A$. Global leastness gives $A\subseteq B$,
so condition~3 gives $\widehat A\ll\widehat B$.
If $D\in\Psi$ is another refined stratum, then $e_A\in D$,
and leastness in $\Psi$ gives $B\subseteq D$ and
$B\in\widehat D$. Hence $B$ is the least carrier.

Suppose compatible figures $\Phi,\Psi$ satisfy
$\Phi\ll\Psi$. In a common admitted parent their cells
are the same as their own cells. A nonempty cell of $\Phi$
contained in a cell of $\Psi$ must be that same parent cell.
Their strata agree, giving $\Phi\subseteq\Psi$.
Finally $A\rho B$ is equivalent to $C_A\subseteq C_B$:
the forward implication uses the carrier of the top cell,
and the reverse uses condition~3 and the least carrier just
proved. Cell containment is transitive, establishing the last
prime carrier axiom.

The two constructions recover the objects canonically.
On the abstract side they identify $x$ with $\Phi_x$ and
recover both orders. On the set side the bijection $e\mapsto A_e$
identifies the canonical refinement ideal of $B$ with its
original set, because $e_A\in B$ exactly when $A\subseteq B$.
The admitted figures are recovered as well. Atelier isomorphisms
transport strata and refinement ideals; bijections of set-atlases
transport cells and their containment. The two recovery maps
commute with these isomorphisms. This proves the stated groupoid
equivalence, including the empty ground set.

\begin{counterexample}[Why the principal subfigure must be admitted]
\label{ex:missing-principal}
On $E=\{a,b\}$ take
\[
 \Phi=\{\{a\},\{a,b\}\},\qquad
 \mathscr H=\{\varnothing,\Phi\}.
\]
This satisfies the five conditions of Theorem~\ref{thm:set-atlas}
with the membership requirement deleted from condition~2 and
all principal refinements interpreted in the ambient
$\operatorname{Fig}(E)$, but it does not recover the prescribed
two ground objects as strata.
\end{counterexample}
\begin{proof}
Bounded unions and consistency of shared principal subfigures
are immediate. The only strict inclusion of strata is
$\{a\}\subset\{a,b\}$, whose principal singleton figure
refines $\Phi$ in the ambient family. For the pair $\Phi,\Phi$,
condition~4 has false premise at the pair $\{a\},\{a\}$:
its principal singleton is not admitted. Pairs involving the
empty figure satisfy the conclusion. The global family is
admissible and separating, with cells $\{a\}$ and $\{b\}$.
Yet $\widehat{\{a\}}=\{\{a\}\}\notin\mathscr H$,
and the incidence poset has only one nonzero irreducible.
\end{proof}

Admissibility also needs an actual least stratum at every point
when infinite families are allowed. On
$E=\mathbb N\cup\{\infty\}$, let
$A_n=\{\infty\}\cup\{m:m\geq n\}$.
Every residual is the nonempty singleton $\{n\}$ and
$A_n\cap A_m=A_{\max(n,m)}$, but $\infty$ has no least
containing member. Thus nonempty residuals and binary-intersection
coverage alone are insufficient. If binary intersections are
unions of common lower strata, adding coverage by the residuals
is sufficient: a residual point of $A$ in another stratum $B$
lies in a common lower stratum; minimality forces that stratum
to be $A$, giving $A\subseteq B$. Every residual point then
has its required least member, and nonempty residuals ensure
that every stratum occurs.

\subsection{Local detection and the original compatibility}
We prove Theorem~\ref{thm:local-relations}. A lieu $p$ has
singleton principal deployment and $U_p=\{p\}$. Every lieu
refining a figure has a carrier in it, so
\begin{equation}\label{eq:lieu-deployment}
 \ell(F)=\bigcup_{x\leq F}U_x.
\end{equation}
For two admitted clique figures,
\begin{equation}\label{eq:figure-disjoint-pairs}
 F\parallel G\quad\Longleftrightarrow\quad
 (\forall x\in F)(\forall y\in G)\ x\delta y.
\end{equation}
The right side excludes a common stratum by irreflexivity and
makes the union an admitted clique. Conversely a compatible
disjoint union has every cross pair distinct and disjoint.
In particular $p\parallel q$ is exactly $p\delta q$ for lieux.

Start with the locally detected strong presentation. If $U_x$
were empty, local detection for $x,x$ would give $x\delta x$,
so condition~1 follows. Substituting local detection into
\eqref{eq:figure-disjoint-pairs}, and then using
\eqref{eq:lieu-deployment}, proves condition~2, even for
arbitrary admitted figures. Distinct compatible strata are
members of one clique; hereditary disjunction along their open
lieux gives condition~3.

For condition~4, let $a\leq x$ but $a\not\leq y$.
Every $p\in U_a$ lies outside $\ell(K)$: otherwise its carrier
$c$ in $K$ would be another internal carrier in $x$, forcing
$a=c\leq y$. Likewise $b\leq y$ but $b\not\leq x$ has
$U_b\subseteq\ell(y)\setminus\ell(K)$.
The hypothesis of condition~4 and local detection give
$a\delta b$. These are exactly the cross pairs not already
contained in one principal clique. Their union is therefore
admitted, proving compatibility of $x,y$.

Conversely retain the two orders and carriers of the atelier
and define $\delta$ by the open-lieu formula with graph
$p\parallel q$. Symmetry is immediate. Nonempty $U_x$ and
$p\not\parallel p$ make it irreflexive. Internal transitivity
gives $U_{x'}\subseteq U_x$ whenever $x'\rho x$, proving
heredity. Condition~3 gives all-pairs principal disjointness.
These are all the strong-system axioms.

We check that this construction recovers the original
compatibility. If $x,y$ were compatible, all distinct strata
in their common deployment satisfy condition~3, so their
union is a $\delta$-clique. Conversely suppose that union
is a clique. Take $p\in\ell(x)\setminus\ell(K)$ and
$q\in\ell(y)\setminus\ell(K)$, with carriers $a,b$ in
$x,y$. Neither carrier lies in $K$, so $a,b$ are distinct
exterior strata. Their $\delta$ relation says $p\parallel q$.
Condition~4 yields the original compatibility of $x,y$.

Compatibility of arbitrary figures is determined on their
strata, so it too is recovered. The original deployments
form a permitted chosen atlas: prime generation gives all
downset subfigures, and a clique union of two admitted figures
has pairwise compatible strata and hence is their original
admitted join. Both constructions thus recover the original
orders and figures. They recover the graph on lieux by its
singleton deployments and the full relation by local detection.
All formulas commute with isomorphisms, proving the stated
equivalence, including the one-figure atelier on the empty
stratum set.

For the final assertion, restricted internal transitivity
already gives $U_x\subseteq U_y$ whenever $x\rho y$.
This makes the induced relation hereditary; nonempty open
lieu sets give irreflexivity, and the stipulated graph test
gives principal disjointness. The carrier presentation is
therefore strong. Theorem~\ref{thm:strong-relations} forces
full internal transitivity. Conversely full transitivity is
necessary for every strong system, so the iff follows.
In the factorization presentation transitivity is already
part of the partial-order hypothesis; the same three checks
give the final assertion.

\subsection{General extension and the limits of uniqueness}
We prove Theorem~\ref{thm:abstract-extension}.
If $H\parallel F$ and $P\ll F$, every stratum $p$ of $P$
has a carrier $c$ in $F$. For $h\in H$, disjointness gives
$h\delta c$, and heredity along $p\rho c$ gives
$h\delta p$. Hence $H\parallel P$, so
\[
 \mathcal N(F)\subseteq\mathcal N(P).
\]

Suppose $P\preceq F$ and $P\subsetneq Q\ll F$.
Choose $x\in Q\setminus P$ and a lieu $p\rho x$.
Every stratum $y\in P$ is distinct from $x$ in the clique $Q$,
so $x\delta y$ and then $p\delta y$. Thus $p\parallel P$.
But $p\ll F$, and a nonzero refining lieu cannot be disjoint
from its target: compatible refinement would imply $p\leq F$,
contradicting zero intersection. This contradicts equal cosupports
and proves incidence maximality.

Conversely, if a maximal $P\ll F$ has different cosupport,
witness absorption gives a nonzero $H\ll F$ with
$H\parallel P$ but $H\not\parallel F$.
Its union with $P$ is a proper incidence enlargement refining
$F$, a contradiction. This proves the equivalence.

Now fix $G\leq F$ and $G'\preceq G$.
Order the figures $Q\ll F$ containing $G'$ by inclusion.
This set is nonempty. A chain has a union which is a downset
clique, hence a figure in the full atlas. Every stratum of the
union still refines $F$, so the union is an upper bound in the
same set. Zorn's lemma gives a maximal member $F'$.
Any larger refinement of $F$ would still contain $G'$, so $F'$
is maximal among all refinements of $F$ and hence a subdivision.
Carrier restriction gives
\[
 G'\leq F'|G\ll G.
\]
Maximality of the subdivision $G'$ forces $F'|G=G'$.
This argument also includes the empty figure.

\begin{proposition}[What the deployment square determines]
\label{prop:deployment-uniqueness}
Let $G$ be a downset in a poset $F$, and let $r:G'\to G$
be monotone. An ordered deployment with inverse image $G'$,
the prescribed isomorphism over $F\setminus G$, and the mixed
incidence rule is unique up to unique isomorphism over those data.
It is the poset
\[
 E=G'\amalg(F\setminus G)
\]
with its two original within-piece orders, no exterior-to-interior
relations, and
\[
 \xi\leq\eta\quad\Longleftrightarrow\quad
 r(\xi)\leq_F\eta
 \qquad(\xi\in G',\ \eta\in F\setminus G).
\]
\end{proposition}
\begin{proof}
The map $c:E\to F$, equal to $r$ on $G'$ and identity on the
exterior, is monotone. Mixed transitivity uses either monotonicity
of $r$ or transitivity in $F$. There is no reverse mixed relation,
so antisymmetry holds. The inverse image of $G$ is exactly $G'$.
Every ordered pair in a proposed deployment is prescribed by the
two within-piece orders or the displayed rule. Consequently the
map fixed by the two identifications is the unique isomorphism.
This proves a statement about ordered deployments.
\end{proof}

\begin{counterexample}[Distinct objects with the same extension diagram]
\label{ex:extension-nonunique}
Even a full strong local atlas with witness absorption and all
compatible-cover gluing can have two different elementary
subdivision extensions realizing the same prescribed deployment
square and the same empty boundary.
\end{counterexample}
\begin{proof}
Take strata $a,b$, identity incidence, refinement $a\ll b$,
and empty primitive disjunction. Every principal deployment is a
singleton, carriers are unique, and internal refinement is the
same two-element partial order. Heredity and principal
disjointness are vacuous. The full figures are $0,a,b$.
The unique lieu is $a$, and $U_a=U_b=\{a\}$, so local
detection holds. The two nonzero figures are incompatible and
both have cosupport $\{0\}$.

Witness absorption is vacuous for nonzero source refinements,
because their cosupports agree. For $P=0$ and $F\ne0$,
the witness $F$ itself works. A compatible incidence cover of a
nonzero figure contains that whole figure, so matching restrictions
uniquely determine the glued refinement. This proves gluing for
arbitrary covers, and the empty case is immediate.

For $F=b$ and $G=G'=0$, both $F'=a$ and $F'=b$ are
subdivisions with the required restriction. Their deployments
are one-point posets and their carrier maps are the unique maps
to the one-point deployment of $b$. Both diagrams have the
prescribed empty inverse image, isomorphic complement and vacuous
mixed rule. Both elementary boundaries are empty. The objects
are nevertheless different.

The strengthened diagram extension also exists for every input
in this model. For $(F,G,G')$ equal to
\[
 (0,0,0),\ (a,0,0),\ (a,a,a),\ (b,0,0),\ (b,b,a),\ (b,b,b),
\]
the extension choices are respectively
$0,a,a,\{a\text{ or }b\},a,b$.
These exhaust the possible inputs. The example leaves open a
separate requirement of a coarsest extension; it disproves
uniqueness merely from the deployment square.
\end{proof}

\subsection{Small support lattices and complement failures}

\begin{proposition}[A ten-element ortholattice]\label{prop:ten-lattice}
Let $I=\{i,i',j,k\}$ and
\[
 T=\{A\subseteq I:(i\in A\Longleftrightarrow i'\in A)\}
       \cup\{\{i\},\{i',j,k\}\}.
\]
Inclusion and set complement make $T$ a ten-element ortholattice
which is not orthomodular. With labels
\[
\begin{array}{c|ccccc}
&0&a'&a&b&c\\\hline
&\varnothing&\{i\}&\{i,i'\}&\{j\}&\{k\}\\[3pt]
&A&A'&B&C&1\\\hline
&\{j,k\}&\{i',j,k\}&\{i,i',k\}&\{i,i',j\}&I,
\end{array}
\]
its covers are
\[
\begin{gathered}
0\prec a'\prec a,\quad 0\prec b,c,\quad a\prec B,C,\\
b\prec A,C,\quad c\prec A,B,\quad A\prec A'\prec1,
\quad B,C\prec1.
\end{gathered}
\]
In particular $c\not\leq C$.
\end{proposition}
\begin{proof}
The paired subsets form an eight-element Boolean cube; the other
two subsets are new. If an intersection of two members belongs
to $T$, it is their meet. Otherwise it contains $i'$ but not $i$
and is a proper subset of $\{i',j,k\}$. Deleting $i'$ gives its
largest allowed subset: no allowed subset of that intersection
can contain $i'$. Thus every meet exists. Complement preserves
$T$ and reverses inclusion, so it gives every join by duality.
Complementary members have meet $0$ and join $1$.

For $x=\{i\}<y=\{i,i'\}$, the intersection with $x^\perp$
is $\{i'\}$ and its allowed meet is empty. Hence
$x\vee(y\wedge x^\perp)=x\ne y$.
The minimal nonzero members are $a',b,c$; the only immediate
successor of $a'$ is $a$. Listing the minimal allowed enlargements
of $a,b,c,A$ gives the displayed covers, and complementation
gives the upper covers. Finally $k\notin\{i,i',j\}$.
\end{proof}

This is the ordinary-powerset reconstruction of the formula in
\cite[part~9, archive sheet~112, PDF p.~113]{Groth}. The handwritten
powerset mark and overwritten drawing do not justify treating every
visible stroke as a covering relation. Reading the mark instead as
``three-element subsets'' yields just three incomparable sets,
not this ortholattice; the mathematical convention is explicit here.

\begin{proposition}[Identity-incidence graph failures]\label{prop:graph-failures}
Give a simple graph identity incidence and internal refinement,
with adjacency as disjunction and all cliques as figures. This
is a strong system with nonempty open-lieu sets and local detection.
For the path $a'\! -\! a\! -\! b\! -\! b'$, its support lattice is
$O_6$, and $F=\{a,b\}$ satisfies
\[
 \cl F\cap N(a)=\{a',b\}\ne\{b\}=\cl(F\setminus\{a\}).
\]
For the five-cycle, singleton supports are singletons, yet the
same relative-complement identity fails. Thus it can fail even
when spatial refinement is reconstructed from singleton supports.
\end{proposition}
\begin{proof}
Principal deployments are singletons, so the carrier and exchange
conditions are identities, and principal disjunction and heredity
are automatic. Every object is a lieu, with $U_x=\{x\}$; all local
tests are exactly the graph tests. Figure refinement is inclusion
and a subdivision is equality: a missing vertex is disjoint from
the smaller clique but not from the larger one containing it.

The path neighborhoods are
\[
 N(a')=\{a\},\quad N(a)=\{a',b\},\quad
 N(b)=\{a,b'\},\quad N(b')=\{b\}.
\]
All intersections of these neighborhoods give exactly the six
closed sets
\[
 \varnothing,\quad E,\quad\{a\},\quad\{b\},\quad\{a',b\},\quad\{a,b'\},
\]
whose order and complements are $O_6$. The adjacent pair has
empty polar and hence full closure, proving its displayed failure.

Number the five-cycle modulo five. The common neighbors of
$i-1$ and $i+1$ consist of $i$, so $\cl\{i\}=\{i\}$.
For $F=\{0,1\}$, its polar is empty and
\[
 \cl F\cap N(0)=\{1,4\}\ne\{1\}=\cl\{1\}.
\]
Inclusion carriers between singleton labelled deployments exist
exactly for identical vertices, so the spatial reconstruction
is identity as asserted. In this edge figure the meet of the
two singleton supports is still zero; the complement failure
alone does not establish a meet failure.
\end{proof}

These identity systems also satisfy the following explicit
witness-absorption property: if $G\ll F$ and a figure is disjoint
from $G$ but not from $F$, some figure refining $F$ has that
property. Indeed $G\subsetneq F$, and a missing vertex of the
clique $F$ is such a witness. For the five-cycle the open support
of each stratum is its singleton, so it is absorbed by that
stratum's refinement ideal as well.

\subsection{External support functions}
For this paragraph an external support assignment on a set $P$
is \emph{normalized} if its supports cover $P$, and \emph{faithful}
if their disjointness detects primitive disjunction. For identity
incidence the open and full external supports coincide. The
stronger \emph{full faithfulness} condition in
\cite[archive sheets~84--85, PDF pp.~85--86]{Groth8} requires
\[
 \bigcup_{z\perp F}|z|^\circ=P\setminus|F|
 \quad\text{for every figure }F.
\]
Here $|F|$ is the union of its stratum supports.

\begin{proposition}[Faithful does not imply fully faithful]
\label{prop:external-faithful}
On two strata $a,b$, take identity incidence, $a\rho b$ besides
identities, and empty disjunction. With $P=\{p,q\}$, the assignment
$|a|=\{p\}$ and $|b|=P$ is normalized, monotone under refinement,
faithful, and separates the points of $P$, but is not fully faithful.
\end{proposition}
\begin{proof}
All deployments are singletons; the proper internal relation
is a partial-order relation and has no nontrivial exchange case.
Thus the data form a strong system, with figures
$\varnothing,\{a\},\{b\}$. All external supports are nonempty and
intersect each other, as required by empty disjunction. They cover
$P$; membership in $|a|$ distinguishes $p$ and $q$.
For $F=\{a\}$ there is no disjoint stratum, so the left side of
full faithfulness is empty and the right side is $\{q\}$.
Notice that inclusion of supports does reflect refinement in
this example; that is a different condition.
\end{proof}

\subsection{Completing labelled spatial deployments}
We prove Theorem~\ref{thm:spatial-completion}. A carrier, when it
exists, is unique: a nonempty label cannot lie inside two distinct
orthogonal labels in the target deployment. Identities and
composites are inclusion carriers. If carriers exist in both
directions, their composite sends each label $S$ to a label of
the same orthogonal family containing $S$. That label must be $S$.
The intermediate inclusion is therefore also equality. Both
carriers fix the actual labels and give identical ordered
labelled deployments, so their source strata coincide by the
injectivity hypothesis. This proves antisymmetry, including for
the relation allowing carriers that are not internal.

For a carrier $h:D_x\to D_y$, let $S_z=h(S_x)$. Order preservation
puts every image below $S_z$. That lower ideal of $D_y$ is precisely
$D_z$, and hence $h$ factors as an internal carrier to $D_z$ followed
by its incidence inclusion. Uniqueness of carriers gives uniqueness
of this factor. Restricting to the principal ideal of a face of
$x$ gives the exchange factor, and all composition compatibilities
follow from uniqueness. The internal relation is reflexive and
transitive, and the preceding antisymmetry argument makes it a
partial order. Pairwise orthogonality in each $D_x$ is precisely
principal disjunction. If $x'\rho x$, then $S_{x'}\subseteq S_x$,
so orthogonality to another label is inherited by $S_{x'}$.
All strong factorization conditions are proved.

For the weaker reading, use the five-cycle with incidence
$0,2<1$ and identities otherwise. Its singleton closed supports
are singleton sets. Every principal labelled deployment is a
family of distinct, disjoint nonempty sets, and all ordered
labelled deployments are different. Both proper incidences
join orthogonal vertices. Nevertheless $0$ and $2$ are not
orthogonal, although both lie in the principal ideal of $1$.
Thus no choice of internal refinement can make these data a
strong system: its principal-disjunction condition already fails.
The orthogonal spatial target in the theorem is essential.

\subsection{An eight-stratum meet failure}
The support calculus tested here is the one formed from open
stratum supports in \cite[archive sheets~92--93, PDF pp.~93--94]{Groth7};
the corresponding absorption condition appears on archive
sheet~86, PDF p.~87. We state the mathematical hypotheses explicitly.
We distinguish primitive disjunction $\delta$ from disjointness
of whole principal figures. Put $x\mathrel{\tau}y$ when $[x]$
and $[y]$ are disjoint figures whose union is a figure. Write
$N_\tau$ and $\cl_\tau$ for its polarity and closure. Define the
open support of a stratum by
\[
 U_x=\cl_\tau([x])\cap N_\tau([x]\setminus\{x\}).
\]
For a fixed figure $F$ and any index subset $A\subseteq F$, put
$S_A=\cl_\tau(\bigcup_{x\in A}U_x)$. The index subset need not
itself be incidence-closed.

\begin{proposition}[Failure within one figure]\label{prop:eight-meet}
Let $M=\{a,b,X,p,e,w,h,g\}$, with incidence $a,b<X$, internal
refinement $p\rho X$, and no other proper generating relations.
Give the seven point strata the disjunction edges
\[
 ab,\ ap,\ bp,\ ae,\ bw,\ pw,\ eh,\ hg,\ gw,
\]
and add $Xa,Xb,Xw$ to primitive disjunction. Use all
incidence-closed primitive cliques as figures. This is a strong
system with nonempty open-lieu sets, local detection, and the
witness-absorption property stated above. Every open support
$U_x$ lies in the refinement ideal of $x$. Nevertheless, in
$F=\{a,b,X\}$,
\[
\begin{aligned}
 S_F\cap N_\tau(S_{\{a\}})&=\{b,p,e\}
       \ne\{b,p\}=S_{\{b,X\}},\\
 S_{\{a,X\}}\cap S_{\{a,b\}}&=\{a,w\}
       \ne\{a\}=S_{\{a\}}.
\end{aligned}
\]
Thus both relative complementation and nonempty-meet preservation
fail for the support calculus of the strata of this one figure.
\end{proposition}
\begin{proof}
The only nonsingleton principal deployment is $\{a,b,X\}$, and
its members are pairwise primitively disjoint. The proper internal
arrow starts at a singleton deployment, so exchange is immediate.
Its heredity requirement holds since $p$ is disjoint from $a,b,w$.
These observations verify the strong conditions. The lieux are
the seven point strata; $X$ has open-lieu set $\{p\}$, while
each point has its singleton. Primitive disjunction of $X$ agrees
with that of $p$, so local detection holds in every case.

For witness absorption, a figure not containing $X$ is a point
clique and every refining figure is its subset. Any missing point
is a witness. A figure containing $X$ must be $F$: no stratum is
disjoint from all of $a,b,X$. A proper figure refining $F$ is
therefore a subset of the point triangle $\{a,b,p\}$. If $a$ or
$b$ is missing, use that point as witness. If the refining figure
is $\{a,b\}$, use $\{p\}$. The remaining case is $\{a,b,p\}$;
no nonempty figure is disjoint from it, since the three point
neighborhoods have empty intersection. The separation premise
then fails. The case of the whole $F$ has a false premise too.

The whole-figure relation $\tau$ has the nine displayed point
edges and makes $X$ isolated. Its point neighborhoods are
\[
\begin{array}{c|cccc}
x&a&b&p&e\\\hline
N_\tau(x)&\{b,p,e\}&\{a,p,w\}&\{a,b,w\}&\{a,h\}\\[3pt]
x&w&h&g&\\\hline
N_\tau(x)&\{b,p,g\}&\{e,g\}&\{w,h\}&
\end{array}
\]
Intersecting neighborhoods shows that every point has singleton
closed support. Since $X$ is isolated, $\cl_\tau([X])=M$ and
$N_\tau(\{a,b\})=\{p\}$. Thus $U_X=\{p\}$, and every other
$U_x=\{x\}$; these are all absorbed by their refinement ideals.
Finally direct intersections of this table give
\[
\begin{gathered}
 S_F=M,\quad S_{\{a\}}=\{a\},\quad
 S_{\{b,X\}}=\{b,p\},\\
 S_{\{a,X\}}=\{a,p,w\},\qquad
 S_{\{a,b\}}=\{a,b,w\}.
\end{gathered}
\]
The two displayed failures follow. This calculation uses $\tau$
and the open supports, throughout; substituting primitive
$\delta$ changes some closed sets.
\end{proof}

\subsection{Proof of the relative-criterion comparison}
We prove Theorem~\ref{thm:criteria}, using the carrier calculus of
Proposition~5.7. Suppose the overlap restrictions in $\mathrm D$
are disjoint. Condition $\mathrm C$ makes $P,Q$ compatible.
Any common stratum would refine both parents and hence their
intersection, so it would belong to both restrictions. This is
impossible, proving $\mathrm C\Rightarrow\mathrm D$.
Taking $K=0$ gives $\mathrm D\Rightarrow\mathrm D_0$.
For distinct compatible elementary $x,y$ and
$p\rho x,q\rho y$, at least one of the two restrictions to
$x\cap y$ is zero. Indeed, if $x,y$ are incomparable, both
are zero; if $x<y$, the internal point of $y$ restricts to zero,
and the other ordered case is symmetric. Thus $\mathrm D$
gives $p\parallel q$, proving $\mathrm L3$.

The $\mathrm L2$ test extends from elementary objects to figures
by the stratum compatibility criterion and
$\ell(F)=\bigcup_{x\in F}\ell(x)$. Refinement decreases these
shadows, so $\mathrm L2$ implies $\mathrm D_0$.
Now assume $\mathrm L2$ and $\mathrm L3$ and take
$p\in\ell(P),q\in\ell(Q)$ under disjoint overlap restrictions.
Let $u=c_F(p),v=c_G(q)$. If $u\ne v$, these are distinct
compatible strata, and $p\rho u,q\rho v$ give disjunction by
$\mathrm L3$. If $u=v$, that carrier lies in $K$, so $p,q$
refine the respective overlap restrictions. Their disjunction
follows from $\mathrm L2$ there. Applying $\mathrm L2$ to all
such pairs proves $\mathrm D$.

For $\mathrm L3+\mathrm{CL}\Rightarrow\mathrm C$, suppose the
overlap restrictions are compatible. It suffices to show that
$x\in P,y\in Q$ are compatible. Put $H=x\cap y$ and take
$p\in\ell(x)\setminus\ell(H)$ and
$q\in\ell(y)\setminus\ell(H)$. Their carriers
$r=c_x(p),s=c_y(q)$ are different, since equality would put
$p$ below $H$. By carrier composition let
\[
 u=c_F(r)=c_F(p),\qquad v=c_G(s)=c_G(q).
\]
If $u\ne v$, they are distinct compatible strata of the parents,
and $\mathrm L3$ gives $p\parallel q$.
If $u=v$, that stratum lies in $K$, so $r\in P|K$ and
$s\in Q|K$. They are therefore compatible, and their distinctness
allows $\mathrm L3$ to give the same conclusion. The exterior
condition $\mathrm{CL}$ now proves $x\bowtie y$; compatibility
by strata proves $P\bowtie Q$.

Under $\mathrm D_0$, disjoint parents have disjoint point shadows.
Conversely, suppose all pairs from $\ell(x),\ell(y)$ are disjoint.
If $x\cap y$ contained a stratum, $\mathrm L1$ would supply a
lieu below it, disjoint from itself, a contradiction. Thus the
intersection is zero. Either $\mathrm{CL}$ or
$\mathrm{CL}_{\rm p}$ now gives compatibility and hence
disjointness. This proves the two remaining arrows to
$\mathrm L2$.

Under $\mathrm L2$ the two absolute antecedents are equivalent.
The $\mathrm{CD2}$ antecedent, tested on singleton lieux outside
$\ell(K)$, implies the exterior-shadow antecedent. Conversely,
if $P|K=0$, no lieu refining $P$ can refine $K$, by the
refinement-infimum property of restriction. The same holds for
$Q$. The exterior-shadow antecedent and figure-level $\mathrm L2$
then give $P\parallel Q$. The added polyhedral qualification is
unchanged. If an alternative polyhedral convention omits the
empty case, that case is already automatic under $\mathrm L2$:
its shadow test directly gives disjointness, hence compatibility.

Equal overlap restrictions are compatible, so $\mathrm C$ gives
the matching condition. For two matching refinements their union
is a figure, refines the union of their parents, and restricts to
the original pieces. Any other gluing is the union of its
restrictions and is equal to this one. Induction, or the
pairwise compatibility criterion together with restriction over
unions, proves finite gluing. Conversely the two-parent cover
gives the matching condition.

For an arbitrary cover, consider the increasing family of glued
refinements over its finite subcovers. Each is a refinement of
the full parent $F$. A stratum of $F$ belongs to some cover
member by primality. Once that member is included, matching
forces the restriction of every later finite gluing to equal
the prescribed local restriction. The stated closure hypothesis
therefore admits the union. It has exactly the required
restrictions and is unique. In a full clique atlas the union is
admitted directly, since every pair of its members is contained
in one of the finite gluings. The restricted-atlas example below
shows why this admission is a separate hypothesis.

\subsection{What the weakened stratum conditions miss}
The next examples concern explicitly stated order and disjunction
conditions. No ambiguous handwritten mark is assigned a definitive
reading by these calculations.

\begin{counterexample}[Least carriers without internal transitivity]
\label{ex:nontransitive}
On $\{p,q,z,u,t\}$, take incidence $z<t$ and refinement generated by
\[
 p\ll z,\quad p\ll u,\quad q\ll u,\quad u\ll t,
 \quad z\ll t,
\]
including identity pairs. Every principal deployment has least
carriers, and every open-lieu set is nonempty. With the two lieux
$p,q$ disjoint, the different strata in every principal deployment
have disjoint open-lieu sets. Nevertheless the induced internal
relation is not transitive, and disjunction defined by the open
lieux is not hereditary under it.
\end{counterexample}
\begin{proof}
All principal ideals are singletons except $[t]=\{z,t\}$.
Into $[t]$, the carriers of $p,z$ are $z$, and those of $q,u,t$
are $t$; every other carrier is forced by a singleton target.
The refinement-minimal objects are $p,q$, and
\[
 U_p=U_z=\{p\},\qquad U_q=U_t=\{q\},\qquad U_u=\{p,q\}.
\]
Thus all the stated point conditions hold, including the sole
nontrivial principal pair $z,t$. But
$p\rho u\rho t$ and $p\not\rho t$. The induced relation has
$p\delta t$ but not $p\delta u$, contradicting heredity along
$u\rho t$. This example is outside the strong-system hypotheses.
\end{proof}

\begin{counterexample}[Comparable pairs are insufficient]
\label{ex:weak-principal}
Partial incidence and internal orders, unique factorization,
exchange and hereditary symmetric irreflexive disjunction do not
force disjunction of every distinct pair in a principal deployment
if the stated disjunction requirement tests only comparable pairs.
The failure persists with nonempty open-lieu sets.
\end{counterexample}
\begin{proof}
First use $y,z<t$, internal identity, and only the edges $yt,zt$.
Internal factorization and exchange are identities; heredity is
automatic. Every comparable distinct pair is disjoint, but
$y,z<t$ are not. If instead the weakened clause requires comparable
pairs to be nondisjoint, the chain $a<t$ with internal identity and
empty disjunction satisfies it and again violates principal
disjunction.

For nonempty open lieux use $\{a,b,c,t\}$, incidence $a,b<t$,
and internal refinement $c\rho t$. The lieux are $a,b,c$,
with edges $ac,bc$ only. Their open sets are singletons, and
$U_t=\{c\}$. The disjunction induced by these open sets has
edges $ac,bc,at,bt$. Unique factors into $t$ are respectively
$a,b,t,t$ for the sources $a,b,c,t$; the only proper internal
arrow has singleton source, so exchange holds. Heredity follows
from the identical open sets of $c,t$. Both comparable pairs
are disjoint, but $a,b$ under $t$ are not. The missing condition
is precisely the all-pairs principal condition.
\end{proof}

\subsection{Four failures inside strong systems}
All four models below use the full family of incidence-closed
primitive cliques. Their internal relations are as stated or
computed from their carriers, so the general strong construction
supplies the prime core and its relative calculus.

\begin{counterexample}[Locality and extension need their hypotheses]
\label{ex:four-strong}
Strong systems alone do not imply $\mathrm L2$ or $\mathrm{CL}$,
even when all open-lieu sets are nonempty. The conditions
$\mathrm L2,\mathrm L3$ and witness absorption together do not
imply either absolute compatibility criterion. Subdivision
extension can also fail in a full strong atlas.
\end{counterexample}
\begin{proof}
For the first assertion take discrete incidence on $p,q,x,y$,
refinement generated by $p\ll x,q\ll y$, and sole primitive edge
$pq$. Principal deployments are singletons, and heredity only
tests the minimal edge endpoints, so this is strong. Both upper
open-lieu sets are nonempty, $U_x=\{p\},U_y=\{q\}$. Their points
are disjoint while $x,y$ are incompatible. Thus $\mathrm L2$
and $\mathrm{CL}$ fail at empty incidence intersection. Any
restricted intersection-class version allowing the empty
intersection also fails here.

A second direct failure of $\mathrm L2$ uses incidence and
refinement equal to the two chains $a<x,b<y$, with primitive
edges $ax,by,ab$. Internal refinement is identity; both principal
pairs are disjoint. The full lieu shadows of $x,y$ are
$\{a\},\{b\}$, which are disjoint, whereas their principal
figures are incompatible. The upper open-lieu sets are empty;
this second example is not asserted under $\mathrm L1$.

For failure of absolute compatibility take incidence and
refinement $a<x,a<y$, with primitive edges $ax,ay$.
Again internal refinement is identity and the system is strong.
Every nonempty figure contains $a$, so no two nonempty figures
are disjoint, and every full lieu shadow is $\{a\}$. This proves
$\mathrm L2$. Distinct compatible strata are $a$ with $x$ or
$y$; their upper open-lieu sets are empty, so $\mathrm L3$ holds.
All nonempty figures have the same cosupport $\{0\}$. Hence
witness absorption is vacuous for a nonzero refinement, and for
a zero refinement the parent is a witness. However the common
incidence figure of $x,y$ is $a$, and both exterior shadows are
empty. The absolute locus tests hold but $x,y$ are incompatible,
refuting $\mathrm{CL}$ and $\mathrm{CL}_{\rm p}$. Every nonzero
refinement of either parent restricts nontrivially to $a$, so the
outside-refinement tests hold as well, refuting $\mathrm{CD2}$
and $\mathrm{CD2}_{\rm p}$.

Finally let incidence and refinement be $a<b<c$, with isolated
$d$, and primitive edges $ab,ac,bc,ad,bd$, omitting $cd$.
Internal refinement is identity, and every principal ideal is
a clique. The complete list of figures is
\[
 0,\quad a,\quad d,\quad a+d,\quad a+b,\quad a+b+d,
 \quad a+b+c.
\]
Here sums denote unions of deployed strata. The figures disjoint
from $a$ and from $a+b$ are exactly $0,d$, so $a$ subdivides the
principal figure of $b$. Only $0$ is disjoint from $a+b+c$;
none of its proper refinements has that cosupport. Thus its only
subdivision is itself, and restriction to $b$ gives $a+b$, not
$a$. The prescribed subdivision of the subfigure does not extend.
This model has empty upper open-lieu sets, so it does not contradict
Theorem~\ref{thm:abstract-extension}.
\end{proof}

\subsection{Gluing does not give relative disjointness}
\begin{counterexample}[All covers can glue while $\mathrm C,\mathrm D$ fail]
\label{ex:gluing-relative}
There is a prime carrier atelier with $\mathrm L1,\mathrm L2$,
$\mathrm{CL}$, $\mathrm D_0$, witness absorption and unique gluing
over every set-indexed incidence cover, but without
$\mathrm C,\mathrm D$ or $\mathrm L3$.
\end{counterexample}
\begin{proof}
Use the four figures $0,\{x\},\{q\},\{x,y\}$.
Incidence is inclusion, and refinement adds
$\{q\}\ll\{x,y\}$. The three nonzero figures are completely
join-prime: a bounded family joining to one of them contains
that figure. Compatibility is equality together with the pair
of elementary strata $x,y$. All bounded joins and generation
are immediate. The least carriers are forced, with $q$ carried
by $y$ and $x$ by $x$ in the principal figure of $y$.
Compatible refinement is incidence. The sole proper internal
arrow is $q\rho y$, so internal transitivity holds. This verifies
all six prime-carrier conditions directly.

The lieux are $x,q$, with $U_y=\{q\}$, so $\mathrm L1$ holds.
No two nonempty figures are disjoint. The full shadows are
$\{x\},\{q\},\{x,q\}$; every two contain a nondisjoint point
pair, proving $\mathrm L2$. The condition $\mathrm D_0$ is
vacuous except when a parent is zero. For the only incompatible
elementary pairs, $x,q$ have nondisjoint exterior points and
$q,y$ have $q$ on both sides; thus their absolute test is false,
and $\mathrm{CL}$ holds. Every nonempty figure has cosupport
$\{0\}$, which proves witness absorption as in the preceding
example.

Every incidence cover of a nonzero figure contains that figure
itself by complete primality. Matching local data are therefore
determined by its chosen refinement, which is the unique gluing.
This holds for arbitrary covers, with repetitions allowed.
But for parents $x\leq y$, take refinements $x\ll x$ and
$q\ll y$. On their overlap $x$, the restrictions are $x$ and
$0$, which are disjoint and compatible. The refinements $x,q$
are neither. Thus $\mathrm C$ and $\mathrm D$ fail.
The compatible distinct strata $x,y$ have nondisjoint open lieux
$x,q$, so $\mathrm L3$ also fails. This is a prime-core example;
it is not claimed to be a strong canonical system.
\end{proof}

\subsection{The missing infinite atlas closure}
\begin{counterexample}[Finite gluing with an omitted stationary union]
\label{ex:restricted-infinite}
A restricted atlas can satisfy the prime core,
$\mathrm L1,\mathrm L2,\mathrm L3,\mathrm{CL}$ and finite gluing,
yet omit a directed union whose restrictions to every stratum
are eventually constant. It can consequently fail arbitrary-cover
gluing.
\end{counterexample}
\begin{proof}
Let $M=\{a_n,x_n:n\in\mathbb N\}$ have discrete incidence,
refinement $a_n\ll x_n$, and primitive disjunction exactly between
different indices. Each principal deployment is a singleton;
heredity preserves the index, so this is strong. The lieux are
$a_n$, and both the full and open lieu shadows of either member
of the $n$th pair are $\{a_n\}$. Distinguishing equal and unequal
indices verifies all four stated local conditions.

A full figure chooses at most one member of each pair. Admit
only figures with finitely many $a_n$, allowing arbitrarily
many $x_n$. This family contains all principal figures and zero,
is closed under subfigures, and is closed under every union
bounded by one of its admitted members. Finite compatible unions
are admitted too. Hence the strong restricted-atlas construction
and Theorem~\ref{thm:criteria} give the prime core and finite gluing.

The figure $F=\{x_n:n\in\mathbb N\}$ is admitted. The increasing
refinements $G_k=\{a_0,\ldots,a_k\}$ restrict to a fixed stratum
$x_n$ as $\{a_n\}$ once $k\geq n$. Their union has infinitely
many $a_n$ and is not admitted; no admitted incidence bound can
contain it. For a direct gluing failure cover $F$ by its singleton
strata and choose $\{a_n\}$ over $x_n$. These data match on their
empty overlaps, but any gluing would contain every $a_n$.
Thus no such gluing exists.
\end{proof}

\subsection{Restricted intersection classes}
We prove Theorem~\ref{thm:sigma}. The proof of
$\mathrm L3+\mathrm{CL}\Rightarrow\mathrm C$ above establishes
its exterior-point condition before using $\mathrm{CL}$.
Apply that argument to the compatible overlap restrictions of
$x,y$: distinct upper carriers give disjunction by $\mathrm L3$,
while coincident upper carriers put the distinct lower carriers
in those compatible restrictions. The restricted biconditional
therefore reduces compatibility exactly to $x\cap y\in\Sigma$.
The carrier/stratum reduction from elementary parents to figures
gives the stated characterization of $\mathrm C$.

For the last sufficient condition put $K=X\cap Y$,
$R=x|K,S=y|K$. Compatibility of the parents forces $K$ to be
empty or elementary. The additional restriction hypothesis makes
$R,S$ empty or elementary too. The common-refinement lemma gives
$x\cap y=R\cap S$. If either restriction is empty this belongs to
$\Sigma$. Otherwise $R,S$ are compatible elementary figures,
and their restricted biconditional puts $R\cap S$ in $\Sigma$.
This proves the positive assertions, including the unrestricted case.

\begin{counterexample}[The restricted criterion with full gluing]
\label{ex:sigma-seven}
The final failure in Theorem~\ref{thm:sigma} occurs with strata
$M=\{a,b,c,p,k,x,x'\}$. The four point strata have singleton
incidence and refinement ideals. The remaining lower ideals are
\[
\begin{array}{c|c|c|c}
 z &[z]&\{w:w\ll z\}&\{w:w\rho z\}\\\hline
 k&\{a,b,k\}&\{a,b,c,k\}&\{c,k\}\\
 x&\{a,b,k,x\}&\{a,b,c,p,k,x,x'\}&\{p,x,x'\}\\
 x'&\{a,b,x'\}&\{a,b,p,x'\}&\{p,x'\}.
\end{array}
\]
Declare exactly the following distinct pairs incompatible:
\[
 (c,k),\ (c,x),\ (p,x),\ (p,x'),\ (k,x'),\ (x,x').
\]
Figures are all incidence-closed compatibility cliques, with
incidence inclusion and refinement tested elementwise. Take
$\Sigma$ to consist of the empty poset and posets with a greatest
element.
\end{counterexample}
\begin{proof}
The incidence and refinement tables are partial orders; their
compatibility relation is hereditary under incidence, as the
listed principal ideals show. Besides the sixteen subsets of
$\{a,b,c,p\}$, the figures are exactly
\[
 \{a,b,k\},\quad\{a,b,k,p\},\quad\{a,b,k,x\},\quad
 \{a,b,x'\},\quad\{a,b,c,x'\}.
\]
Indeed a figure containing $x$ cannot contain $c,p,x'$; one
containing $k$ cannot contain $c,x'$; and one containing $x'$
cannot contain $p,k,x$. This proves the complete count.

Bounded joins are unions, principal figures are completely
join-prime, and compatibility is detected by strata. Lifted
refinement is a partial order: within a compatible figure,
refinement-comparable strata are already incidence-comparable,
by the table. Its maximal refinement elements therefore determine
that incidence downset. Mutual lifted refinement identifies
these maximal elements in the finite stratum poset, and hence
identifies the figures.

All figure carriers can be listed. A source $a$ or $b$ always
has itself as least carrier. A source $c$ is carried by $c$
when that point occurs, and otherwise by $k$ whenever it refines
the target. A source $p$ is carried by whichever of $p,x,x'$
occurs; these possibilities cannot coexist in one figure.
A source $k$ is carried by $k$, a source $x$ by $x$, and a
source $x'$ by the unique occurring member of $\{x',x\}$.
The proper internal arrows are $c\rho k,p\rho x',x'\rho x$
and $p\rho x$, so transitivity holds. Each joins incompatible
strata; the remaining refinement comparisons are incidences.
Thus compatible refinement is incidence, and all six prime-core
conditions hold.

Give the strata the ground supports
\[
\begin{gathered}
 A_a=\{a\},\quad A_b=\{b\},\quad A_c=\{c\},\quad A_p=\{p\},\\
 A_k=\{a,b,c\},\qquad A_x=\{a,b,c,p\},\qquad A_{x'}=\{a,b,p\}.
\end{gathered}
\]
These are their full lieu shadows. Open-lieu sets are the point
singletons and $U_k=\{c\}$, $U_x=U_{x'}=\{p\}$.
Thus $\mathrm L1$ holds. Distinct compatible strata have different
open points, proving $\mathrm L3$.
For a figure put $A_F=\bigcup_{z\in F}A_z$.
Whole-figure disjointness is equivalent to $A_F\cap A_G=\varnothing$.
In one direction, disjoint supports exclude every forbidden pair
and any common stratum. In the other, compatible elementary
strata with overlapping supports are incidence-comparable by
the table; overlapping supports of two compatible figures
therefore force a common smaller stratum. This proves the
equivalence and $\mathrm L2$.

For compatible elementary pairs, their common deployment is
empty or principal and their exterior support sets are disjoint.
For the incompatible pairs $ck,cx,px,px'$, the common deployment
is empty but the exterior supports intersect. For $xx'$ the
common deployment is $\{a,b\}$ and the exterior supports still
share $p$. For $kx'$ the exterior supports are $\{c\}$ and
$\{p\}$, but their common deployment $\{a,b\}$ has no greatest
element. These cases prove the exact restricted biconditional.

We next prove $\mathrm D$ for all figures. For compatible
parents, the preceding support argument gives
$A_{F\cap G}=A_F\cap A_G$. Moreover
\[
 A_P\subseteq A_F\quad(P\ll F),\qquad
 A_{P|K}=A_P\cap A_K\quad(K\leq F).
\]
For the second equality, every ground point of a figure belongs
to the open-lieu set of a unique stratum. Carriers preserve that
open point, by the internal-arrow table. If a point is in both
$A_P$ and $A_K$, the corresponding stratum of $P$ has its carrier
in $K$, by uniqueness of the open-point stratum in $F$. It
therefore belongs to $P|K$. The other inclusion is immediate.
Disjoint overlap restrictions now force $A_P\cap A_Q$ to be
empty, since that intersection lies in $A_K$; hence $P\parallel Q$.

Matching refinements also glue. Their union is an incidence
downset; it suffices to exclude every forbidden cross-pair.
Except for $kx'$, each forbidden pair has the same open point,
or reduces to such a pair: $cx$ reduces to $ck$ because
$k\leq x$. Suppose such a pair occurred in opposite refinements
$P,Q$ of compatible parents $F,G$. Its carriers have the same
open point. Since they are compatible strata and $\mathrm L3$
holds, those carriers must be equal and lie in $K=F\cap G$.
Matching restrictions would then put the forbidden pair in
one figure $P|K=Q|K$, a contradiction.
For the remaining pair, suppose $k\in P,x'\in Q$.
The parent $F$ contains $k$, since the only possible targets
of $k$ are $k,x$. A target of $x'$ in $G$ is $x'$ or $x$;
compatibility with $F$ excludes the first, so $G$ contains $x$
and hence $k$. Thus $k\in K$ and $k\in P|K$. Equality of
restrictions puts $k$ in $Q$, contradicting $x'\in Q$.
The union is therefore admitted and is the unique gluing.
Theorem~\ref{thm:criteria} gives finite-cover gluing. Arbitrary
covers in this finite atlas reduce to finitely many distinct
members, with repetitions forced to have identical local data.

Finally use parents $k\leq x$ and refinements $k\ll k,x'\ll x$.
Their restrictions to $k$ are $k$ and $\{a,b\}$, which are
compatible. But $k,x'$ are incompatible, precisely because their
common deployment has two incomparable maximal elements.
Thus $\mathrm C$ fails with every asserted extra condition intact.
\end{proof}


\section{Finite figures and binary event structures}\label{sec:events}
We prove Theorem~\ref{thm:events}. Deployment is order reflecting because figures are joins of their deployed irreducibles. Each principal ideal therefore injects into the powerset of the finite set $\delta(F)$ and is finite. Bounded suprema make it a lattice. In a distributive such ideal, $m\le G\vee H$ implies
\[
 m=(m\wedge G)\vee(m\wedge H),
\]
so irreducibility makes $m\le G$ or $m\le H$. Induction and finiteness of the ideal give primality for every bounded family.

Primality gives $\delta(\bigvee F_i)=\bigcup\delta(F_i)$ for bounded families. Let $A\subset M$ be a finite downward-closed compatibility clique. Join its elements successively. The partial join is compatible with the next element because every deployed member is, using downward heredity of compatibility and the stipulated pairwise criterion. Its join exists, and primality gives $\delta(\bigvee A)=A$. The empty set gives zero. Thus deployment is onto the asserted configuration set. Conversely the ideal below a configuration is the downset lattice of its finite causal poset, and is distributive.

Failure of compatibility is a symmetric irreflexive upward-hereditary conflict relation, and each causal ideal $\delta(m)$ is finite. Conversely finite configurations of such an event structure satisfy all the stated hypotheses; their irreducibles are exactly the principal causes $\downarrow e$. These constructions are inverse on objects and isomorphisms, proving the theorem.

Irreducibility alone is insufficient. Take the lattice $0<d<a,b,c<1$ with $a,b,c$ incomparable. Its irreducibles are $d,a,b,c$, every pair is compatible, and every figure is their join. But $\{d,a,b\}$ is a downward-closed clique which is not any deployment, since $a\vee b=1$ also deploys $c$. Moreover general finite-consistency event structures are broader than binary conflict: three independent events with every proper subset consistent but the triple inconsistent cannot be encoded by pairwise conflict.

\subsection{Completion and whole-figure supports}
Let $E$ be the event poset obtained in Theorem~\ref{thm:events}.
Thus conflict $\#$ is symmetric, irreflexive and hereditary upward,
and every principal past $\downarrow e$ is finite. A configuration
is a downward-closed conflict-free subset. Write $\mathcal C_f(E)$
for finite configurations and $\mathcal C(E)$ for all configurations.
The orders are inclusion, with empty configuration $0$.

For two configurations define whole-figure disjointness by
\[
 F\perp G\quad\Longleftrightarrow\quad
 F\cap G=\varnothing\ \text{ and }\ F\cup G\in\mathcal C(E).
\]
On events use the relation
\[
 e\perp_E f\quad\Longleftrightarrow\quad
 \neg(e\# f)\ \text{ and }\
 \downarrow e\cap\downarrow f=\varnothing.
\]
These relations concern whole figures. They are different from the
primitive disjunction between open cells used in Section~7.

\begin{proposition}[Configuration completion and supports]
\label{prop:event-domain}
There is a canonical order isomorphism
\[
 \mathcal C(E)\simeq\operatorname{Idl}(\mathcal C_f(E)),
\]
where an ideal is a nonempty directed downward-closed subset.
The compact elements of $\mathcal C(E)$ are its finite configurations,
and its nonzero completely join-prime elements are exactly the
principal pasts $\downarrow e$. Every pairwise bounded family has
a join. Moreover there are canonical complete ortholattice
isomorphisms
\[
 \Sigma(E,\perp_E)\simeq
 \Sigma(\mathcal C_f(E)\setminus\{0\},\perp)
 \simeq\Sigma(\mathcal C(E)\setminus\{0\},\perp).
\]
If $0$ is retained as a universally self-orthogonal vertex,
the same support lattice results, with bottom support $\{0\}$.
\end{proposition}
The configuration-domain assertions are the specialization of
Nielsen--Plotkin--Winskel's coherent prime-algebraic domain
correspondence \cite[Definitions~17--18 and Theorems~8--9, p.~102]{NPW}.
We give the elementary argument and the support comparison with
all empty and infinite cases explicit.

\begin{proof}
Every principal past is conflict-free: a conflict between two of
its elements would propagate to a self-conflict at its maximum.
A family of configurations is bounded exactly when its union is
conflict-free. The union is then its join. Pairwise boundedness
already excludes every possible conflict, proving the coherence
assertion and the formula for directed joins.

For a configuration $C$, its finite subconfigurations form an ideal.
Their union is $C$, since each event belongs to its finite principal
past. Conversely the union of an ideal is a configuration: any
two events lie together in a common ideal member. A finite
subconfiguration of that union lies in one ideal member by
directedness and therefore belongs to the ideal. These operations
are inverse, including when $E$ is empty.

A finite configuration contained in a directed union is contained
in one member, by finitely many witnesses and directedness. It is
therefore compact. Every configuration is the directed union of
the finite unions of its principal pasts, so a compact configuration
must itself be finite. If $\downarrow e\subseteq\bigvee_i C_i$,
the join is a union and $e\in C_i$ for some $i$. Downward closure
gives $\downarrow e\subseteq C_i$. Conversely a nonzero completely
join-prime configuration, written as the union of its principal
pasts, must equal one of them. This proves the domain assertions.

The event orthogonality is hereditary downward in both variables.
For configurations, finite or infinite, one has
\begin{equation}\label{eq:event-pairwise}
 F\perp G\quad\Longleftrightarrow\quad
 (\forall e\in F)(\forall f\in G)\ e\perp_E f.
\end{equation}
In the forward direction causal closure puts the two pasts in
the two disjoint configurations. Conversely the displayed condition
excludes cross-conflicts and any common event, proving the claim.

Let $V$ be either nonempty configuration set in the statement.
For $\mathcal U\subseteq V$ and $S\subseteq E$, put
\[
 U(\mathcal U)=\bigcup_{F\in\mathcal U}F,
 \qquad J(S)=\{F\in V:F\subseteq S\}.
\]
Equation~\eqref{eq:event-pairwise} gives
\[
 N_V(\mathcal U)=J\bigl(N_E(U(\mathcal U))\bigr).
\]
Every event polar is downward closed. If $S$ is downward closed,
then $U(J(S))=S$, because the finite past $\downarrow e$
belongs to $J(S)$ for every $e\in S$. Consequently, for a closed
event support $S$,
\[
 N_V(J(S))=J(N_E(S)),\qquad N_V^2(J(S))=J(S).
\]
Conversely, applying the same formulas twice to a closed
configuration support $\mathcal U$ yields
$\mathcal U=J(N_E^2(U(\mathcal U)))$ and
$U(\mathcal U)=N_E^2(U(\mathcal U))$.
Thus $J$ and $U$ are inverse order and orthocomplement isomorphisms.
This proves the assertion for both finite and arbitrary configurations.
The test objects used here are principal pasts; a singleton event
need not be a configuration.

Finally, if $0$ is retained, every polar contains it because it
is orthogonal to every figure, including itself. Adding or removing
this one vertex gives the asserted support isomorphism.
\end{proof}

\subsection{Which data a configuration domain forgets}
Incidence and compatibility are recovered from configurations.
Refinement requires additional structure. For clarity the
translation of the refinement-by-strata rule is
\[
 G\ll F\quad\Longleftrightarrow\quad
 (\forall e\in G)(\exists f\in F)\ e\ll f,
\]
where refinement of events means refinement of their principal
configurations. If subdivision is defined to preserve whole-figure
cosupport, its translation is
\[
 G\preceq F\quad\Longleftrightarrow\quad
 G\ll F\ \text{ and }\ N(G)=N(F).
\]
Here $N(F)$ denotes all configurations disjoint from $F$.
Triple-polar equality makes this equivalent to equality of the
double-polar supports. Neither formula supplies the additional
refinement relation from the event structure alone.

The following examples use a precise elementary convention. The
principal deployment of a stratum $f$ is $\downarrow f$. A carrier
of $e\ll f$ is the least $d\leq f$ with $e\ll d$; internal
refinement $e\rho f$ means that this carrier is $f$ itself.
A primitive disjunction is symmetric and irreflexive, is hereditary
under internal refinement in both variables, and relates every
two distinct members of one principal deployment. Figures in the
examples are exactly its downward-closed cliques.

For an explicit example, take two incomparable conflicting events
$a,b$. The configurations are $0,A=\{a\},B=\{b\}$.
Use incidence alone as refinement in the first system, and add
$A\ll B$ in the second. Both are partial orders; their nonzero
principal deployments are singletons, so carriers exist uniquely,
and compatible refinement is incidence. Internal refinement is
transitive. Set primitive disjunction on the two strata to be empty;
its heredity and the distinct-strata condition within one principal
deployment are vacuous. The first system has two minimal refinement
objects, or lieux, and the second has only $a$. All nonzero figures
have polar $\{0\}$, so $A$ subdivides $B$ only in the second system.

Primitive internal disjunction can also vary while the event
structure stays fixed. Take two chains $a<a'$ and $b<b'$, with
every cross-chain pair in conflict. The configurations are
\[
 0,\quad\{a\},\quad\{a,a'\},\quad\{b\},\quad\{b,b'\}.
\]
Set refinement equal to incidence. One primitive relation has
exactly the unordered pairs $\{a,a'\}$ and $\{b,b'\}$; the other
also has $\{a',b'\}$. Internal refinement is identity, so heredity
holds. In each case the downward-closed primitive cliques are
exactly the five displayed configurations: any mixed downset
contains the missing pair $\{a,b\}$. Thus the causal order,
conflict, configurations and whole-figure supports agree, while
the primitive relation differs.

This also distinguishes a cosupport definition from a maximality
theorem. A witness-absorption condition saying that a figure
disjoint from $P\ll F$ but not from $F$ can be chosen to refine
$F$ does not alone make subdivisions incidence-maximal. In the
last example it holds: for nonzero $P$ its premise is empty, and
for $P=0$ one can use $F$ itself. Yet $\{a\}$ is a proper
cosupport-preserving incidence subfigure of $\{a,a'\}$.
The upper strata in this example have no open lieux. Under the
strong carrier hypotheses, nonempty open-lieu sets and witness
absorption give the exact maximality conclusion of
Theorem~\ref{thm:abstract-extension}; the full atlas then also
gives subdivision extension.


\section{A definable cell atelier}\label{sec:cellatelier}
The cell figures, their two orders, and primitive disjunction are those defined in Section~1. We first establish the carrier calculus needed for the three cell theorems stated there.

\begin{proposition}[Carriers and finite figure calculus]\label{prop:cellcalculus}
Incidence and refinement are partial orders, and incidence implies refinement. A refinement has a unique monotone carrier
\[
 c_{G,F}:G\longrightarrow F,\qquad C_B\subseteq C_{c_{G,F}(B)}.
\]
Two figures are compatible precisely when their distinct elementary subfigures have pairwise disjoint top cells. Their join is then the union of their elementary subfigures. The principal incidence ideal below $F$ is the finite distributive lattice of downsets of its closed-stratum poset. Elementary figures are precisely its nonzero join-irreducibles and are prime for bounded joins.

If $G\ll F$ and $H\le F$, the induced restriction is
\[
 G|H=c_{G,F}^{-1}(H),\qquad |G|H|=|G|\cap|H|.
\]
These restrictions compose. Compatible figures have empty common subfigure if and only if their supports are disjoint.
\end{proposition}
\begin{proof}
Reflexivity and transitivity of refinement are immediate. If $F\ll G\ll F$, each $F$-cell lies in a $G$-cell which lies in an $F$-cell. Disjointness of the $F$-cells makes the first and last equal, and hence all three equal. Thus $F=G$. Incidence is a partial order and preserves open cells, proving its inclusion in refinement.

Carriers are unique because target open cells are disjoint. If $B'\subseteq B$, then $C_{B'}\subseteq B\subseteq c(B)$. The target cell $C_{c(B')}$ meets the closed target stratum $c(B)$, and the frontier property gives $c(B')\subseteq c(B)$. This proves monotonicity. The inverse image of a downset is a downset. A point belongs to $|H|$ exactly when its target cell belongs to $H$, which proves the support formula. Uniqueness of carriers proves transitivity of restriction.

Suppose the elementary subfigures under consideration have disjoint top cells whenever they are distinct. Unite those cells and their boundary data, identifying equal elementary figures. Every old closed stratum remains the union of its old downward family. Its frontier is unchanged, so this is a finite frontier partition with closed support. Distinct resulting strata cannot have the same closed support: their top cells would be disjoint relatively open dense subsets of that support. The union is therefore a figure containing both given figures. Conversely a common upper figure provides exactly such disjoint open cells and compatible boundary data. The union is its least subfigure containing both, proving the join assertion.

The description of principal ideals as downsets proves distributivity, the elementary description of irreducibles, and bounded-join primality. If compatible figures have a point in both supports, the unique cell of their join containing that point belongs to both, so its elementary figure is a nonempty common subfigure. Conversely a common nonempty subfigure supplies a common point. Disjoint supports permit their union as a figure, since closure is taken in the fixed ambient $X$. This proves the last assertion.
\end{proof}

The carrier of an elementary refinement has a useful minimality property: the least target subfigure refined by it is the elementary figure below its top-cell carrier. Indeed all lower carriers lie below that carrier, and the carrier itself must occur in any target subfigure containing the refinement. In particular, refinement is detected on elementary figures. A strict incidence chain has length at most the dimension of its largest stratum plus one, hence at most $n+1$.

\begin{proof}[Proof of Theorem~\ref{thm:extension}]
Keep each old open cell outside $|G|$ and replace the cells of $G$ by those of $G'$. They form a finite partition of $|F|$. A stratum outside $G$ has the same closed support as before. Its intersection with $|G|$ is the union of the old boundary cells below it; the carrier restriction partitions that intersection by new cells of $G'$. Hence its closure is again a union of whole cells. A new cell has closure inside $|G|$ because this subset is closed in $X$, and the frontier property there is the one for $G'$. Thus the new partition is a figure, with the stated closed strata and support. It refines $F$ and restricts to $G'$.

If $B\subseteq A$, the target cell containing $C_B$ meets $A$, so the old frontier property puts its closure $c(B)$ below $A$. Conversely $B\subseteq c(B)$, so $c(B)\subseteq A$ implies $B\subseteq A$. This proves the mixed incidence formula.

For any other $H$ as stated, its cells over $G$ are exactly those of $G'$. A cell whose carrier lies outside $G$ lies in the corresponding old open cell, which is an $E$-cell. Therefore every $H$-cell lies in an $E$-cell and $H\ll E$. When $F$ is elementary and $G$ its boundary, every new stratum is below the same top stratum, giving the last assertion.
\end{proof}
This theorem does not say that carriers reflect every order relation. Subdivide $[0,1]$ at $1/2$. Its two open half-intervals have the same old interval carrier, while their closed supports are distinct and incomparable. Order reflection holds in the stated mixed formula, with the second stratum outside the replaced boundary.

\begin{proof}[Proof of Theorem~\ref{thm:moderation}]
We use compatible definable triangulation: a closed bounded definable set $A$, with finitely many specified definable subsets, has a definable homeomorphism $h:|K|\to A$ from a finite simplicial complex such that each specified subset is a union of images of open simplices \cite[Theorem~4.4, p.~42, proof pp.~42--44]{Coste}. This theorem holds in any fixed o-minimal expansion of a real closed field.

The coordinate map $b(t)_i=t_i/(1+|t_i|)$ is a definable homeomorphism $R^n\to(-1,1)^n$, with inverse $s_i/(1-|s_i|)$. Put $A=\overline{b(X)}^{\,R^n}$. Triangulate $A$ compatibly with $b(X)$, $b(S)$, and the images of all cells in the given finite collection. For each open simplex $\sigma$ whose image lies in $b(S)$, use the cell $b^{-1}(h(\sigma))$. An open simplex is semialgebraically homeomorphic to an open box of its dimension, so these are generalized cells partitioning $S$.

The closure of $h(\sigma)$ in $b(X)$ is $h(\overline\sigma)\cap b(X)$. Compatibility makes each open face either wholly inside or wholly outside $b(X)$. Since $b(S)$ is closed in $b(X)$, all faces in $b(X)$ occurring in this closure lie in $b(S)$. Thus the selected cells satisfy the frontier condition with relative closures in $X$. Compatibility with all old cells makes each new cell lie in one old cell, as required. The empty support uses the empty figure.
\end{proof}
Only a finite compatibility triangulation is used here. There is no assertion that a prescribed triangulating map can be subdivided linearly, or that infinitely many refinements have a common finite refinement. Indeed the subdivisions of $[0,1]$ at every $t\in(0,1)$ cannot have such a common refinement: every singleton $t$ would have to be a cell.

\begin{proof}[Proof of Theorem~\ref{thm:cellpowerset}]
For any family $\mathcal U\subseteq M_X$, let $T=\bigcup_{m\in\mathcal U}m^\circ$. The definition of primitive disjunction gives
\[
 N(\mathcal U)=S_{X\setminus T}.
\]
Every point $x\in X$ defines a singleton elementary figure, since parameters are allowed. Thus the union of the top cells of $S_Q$ equals exactly $Q$, even for nondefinable $Q$. Applying the formula twice gives $N(S_Q)=S_{X\setminus Q}$ and $N^2(\mathcal U)=S_T$. Consequently the closed subsets are exactly the $S_Q$, and their order and complements are the powerset order and complement. The same formula shows that joins and arbitrary meets are preserved.

For a subfigure, its open strata partition its support. Its elementary figures have exactly these top cells, so both descriptions give the same $S_{|F|}$. Finally $\{x\}\ll F$ if and only if $x\in|F|$. Equality or inclusion of supports is therefore detected by their point refinements. Singleton figures are precisely the minimal nonempty figures for refinement: every nonempty figure admits a point refinement, and a singleton admits no proper nonempty refinement.
\end{proof}

\begin{proposition}[The two connectedness statements]\label{prop:cellcomponents}
The connected components of the undirected refinement graph of $M_X$ are exactly the definably connected components of $X$. For each finite figure, the connected components of its incidence graph are exactly the definably connected components of its support.
\end{proposition}
\begin{proof}
Every generalized cell is definably connected, and so is its closure in $X$. For the latter assertion, a definable clopen separation would restrict to a separation of the dense cell, and density rules out an unused nonempty open part. Hence every elementary support lies in one definably connected component of $X$. Refinement implies support inclusion, so this component is constant along refinement edges. Every point of an elementary support also refines that elementary figure.

Conversely two points in the same definably connected component can be joined by a continuous definable path \cite[Theorem~3.9, p.~28]{Coste}. Its image is closed and bounded: coordinate extrema give boundedness, and the minimum of distance from a proposed missing closure point gives closedness \cite[Theorem~3.4, pp.~26--27]{Coste}. Triangulate the image compatibly with the two endpoints. Its finite incidence graph is connected; otherwise a component and its complement would be disjoint finite unions of closed strata, hence a nontrivial definable clopen partition of the connected path image. Incidence edges are refinement edges, which proves the first assertion.

For a finite figure, each incidence component is relatively clopen by the same finite-union argument. Each closed stratum is connected and incident strata meet. A connected finite incidence graph therefore has connected union, for example by successively adjoining strata along a spanning tree. Thus its incidence components are exactly its support components.
\end{proof}

These conclusions distinguish the two orders. On $X=R$, the one-cell figure $R$ is isolated in the global incidence graph, whereas the singleton $\{0\}$ is a different incidence component and is incompatible with it. They are connected by refinement. There is also no general polyhedral intersection condition: on a circle partitioned into two vertices and two open arcs, the two closed arcs meet in a two-point figure, which is neither empty nor elementary. With one vertex and its complementary open arc, the top stratum has only one immediate predecessor. Finally, replacing fixed ambient closures by closures relative to arbitrary nonclosed supports would spoil the calculus: the one-cell figure $(0,1)$ and the point $1$ have disjoint supports, yet cannot both be closed subfigures of a common figure on $(0,1]$. The relative-closed-support hypothesis in the definition prevents this failure.

\section{Oriented maquettes}
We first justify the rank and diamond assertions in
Theorem~\ref{thm:cw}. A finite simplicial complex homeomorphic
 to an $m$-sphere is pure of dimension $m$: the interior of a
maximal simplex has a neighborhood open in that simplex alone,
so invariance of dimension fixes its dimension. Consequently
all maximal chains below $x$ have $d(x)$ elements. Appending
$y$ to a maximal chain below it, when $y\lessdot x$, gives a
maximal chain below $x$, so $d(x)=d(y)+1$. Saturating finite
intervals proves that this is a rank function.

For a rank-two interval $[z,x]$, a maximal chain ending at $z$
is a codimension-one simplex in the sphere $|\Delta(P_{<x})|$.
Its link is a finite discrete set of $k$ points. Local top
homology at an interior point of that simplex is
$\widetilde H_0$ of the link, of rank $k-1$; in a topological
sphere it is $\Z$. Thus $k=2$, giving exactly two intermediate
cells. An edge has lower sphere $S^0$, hence exactly two vertices.
This argument uses topological spheres and does not impose a PL
condition.

For the realization use Bj\"orner's spherical-interval construction. Give $|\Delta(P)|$ the closed cells $|\Delta(P_{\le x})|$, each the cone on the prescribed sphere $|\Delta(P_{<x})|$. The corresponding open cell is the difference, an open ball. Each point has a finite supporting chain with a largest element, so these open cells partition the realization. Finite principal ideals make cell closures finite subcomplexes. The cell weak topology equals the simplicial weak topology because every simplex lies in the closed cell of its largest vertex. The continuous barycentric coordinate functions separate points, so the realization is Hausdorff. Thus this is a regular CW complex, without any global finiteness restriction. If $P$ is empty the construction gives the empty complex.

Choose orientations by increasing dimension. Vertices are prescribed. For an edge, the reduced zero-cycle group of its two endpoints is generated by their difference; the opposite-coefficient condition chooses exactly one of its two orientations. For a cell $x$ of dimension at least two, suppose the lower cells have been oriented. The prescribed chain
\[
 c_x=\sum_{y\lessdot x}\epsilon(y,x)y
\]
is a top-dimensional cycle on the boundary sphere: its boundary coefficients cancel in signed diamonds. The top cycle group of that sphere is $\Z$, and a fundamental cycle has coefficient $\pm1$ at every top cell. Since $c_x$ does too, it is one of its two generators. Choose the orientation of $x$ accordingly. This constructs all the required signs. Conversely regular cell boundaries are spheres and cellular boundaries satisfy $\partial^2=0$; rank-two intervals have two intermediate cells. These give all the stated conditions, proving the theorem.

\subsection{Explicit incidence laws in low dimensions}
The distinction between oriented elements and orientation orbits is
already necessary for consistency. If an irreflexive relation on a
nonempty set with free reversal is invariant under reversing either
argument, it cannot be ordinary inequality. Indeed
$x\ne\bar x$ would give $x\perp\bar x$, hence
$\bar x\perp\bar x$. The relation
$x\perp y\Longleftrightarrow q(x)\ne q(y)$ is a consistent
alternative. It distinguishes underlying cells, not their two
orientations.

We state the low-dimensional incidence conditions explicitly,
writing $\triangleleft$ for immediate incidence.
Let the orbit poset be graded in nonnegative dimensions, let immediate
incidences project exactly to covers, and require simultaneous
reversal. Choose positive vertex orientations. Impose these two laws:
\begin{enumerate}
\item Each oriented edge has exactly one positive and exactly one
negative vertex immediately below it.
\item If $z<x$ has dimension difference two, then for each orientation
$Z$ of $z$ and each orientation $X$ of $x$ there is exactly one
oriented intermediate cell $Y$ with
$Z\mathrel{\triangleleft}Y\mathrel{\triangleleft}X$.
\end{enumerate}
Unlike Theorem~\ref{thm:cw}, these two laws do not prohibit both
orientations of a facet from being incident to the same oriented cell.
That distinction is essential in dimension two.

\begin{proposition}[The graph and boundary dictionary]\label{prop:lowdim-dictionary}
In dimensions at most one, the displayed incidence data are exactly
graphs specified by a vertex set $V$, an oriented-edge set $E$ with
free reversal, and an origin map $s:E\to V$. The terminus is
$t(e)=s(\bar e)$. Loops and parallel edges are allowed.

In dimension at most two, the oriented edges below each chosen
two-cell orientation form a disjoint union of finite directed cycles
and two-sided infinite directed paths. Every incident vertex occurs
once as origin and once as terminus. Conversely any such boundary
data determine incidence data satisfying the two laws. If each cell
has finitely many faces, only finite cycles occur.
\end{proposition}
\begin{proof}
The first law gives the unique origin $s(e)$. Reversal gives the
unique negative terminal vertex $t(e)=s(\bar e)$. Thus the signed
incidences are exactly
\[
 s(e)^+\triangleleft e,\qquad t(e)^-\triangleleft e,
\]
and their simultaneous reversals. These formulas also reconstruct
the incidence relation from any graph of the stated kind.

Fix a two-cell orientation $f$ and let $B_f$ be its incident
oriented-edge set. Applying the second law to the positive and
negative orientations of a boundary vertex gives exactly one
outgoing and one incoming edge in $B_f$. Following the outgoing
edge is therefore a bijection of the boundary vertex set.
Its orbits are finite cycles or free $\Z$-orbits, the latter
being two-sided infinite paths. Conversely these conditions give
exactly one two-step path for each endpoint orientation, so they
recover the second law. A finite face set excludes an infinite orbit.
\end{proof}
For example, the edges $e_i:v_i\to v_{i+1}$, $i\in\Z$, may
all be placed below one formal positive two-cell. Every incidence
chain has length at most two, yet its boundary is a line.
Finite dimension does not imply finite cell closure.

There is also a substantive difference between matching and
unmatched signs in an incidence-lifting condition. Write $D_X$ for
the immediate predecessor set of $X$. A condition that, under
identity refinement, transfers an incidence into either $X$ or
$\bar X$ to an incidence into $X$ forces
$D_X=D_{\bar X}$. With the displayed laws this leaves exactly
graphs consisting of loops and isolated vertices. To see this,
simultaneous reversal makes every $D_X$ reversal-stable. Any
two-step chain then has both $Y$ and $\bar Y$ as intermediates,
contrary to uniqueness. Thus no cell of dimension at least two
remains. At an edge, reversal stability and the unique endpoint
law make its positive origin and negative terminus lie over the
same vertex. Conversely loop graphs satisfy all these conditions.
If the conclusion uses the same upper sign as the hypothesis,
the lifting condition under identity refinement is tautological.

\begin{proposition}[The exact regular class in dimension two]
\label{prop:lowdim-regular}
For the preceding incidence data in dimensions at most two, assume
each cell has finitely many faces. They arise from a regular CW
complex exactly when no edge is a loop, every two-cell boundary
is nonempty and connected, and no such boundary uses both
orientations of one underlying edge. No global bound on the
number of cells or on vertex valence is required.
\end{proposition}
\begin{proof}
A finite connected directed boundary is one cycle, by
Proposition~\ref{prop:lowdim-dictionary}. Its vertex set consists
of the distinct elements in one orbit. The two exclusions make
its underlying edges distinct and rule out a one-edge loop.
A two-edge cycle with distinct parallel edges is allowed.
Thus each prescribed boundary embeds as a circle subcomplex of
the underlying graph. Realize non-loop edges by closed intervals
and attach a disk to each boundary circle by a homeomorphism.
The resulting closure-finite weak cell structure is a regular
CW complex. Conversely a regular edge is an embedded closed
interval and a regular two-cell has an embedded boundary circle.
It has neither a loop edge nor a repeated underlying edge, and
its boundary is connected. These give precisely the conditions.
\end{proof}

Excluding loops alone is insufficient. Take two vertices $a,b$,
one non-loop edge $e$ from $a$ to $b$, and one formal two-cell
whose positive boundary contains both $e$ and $\bar e$.
Reverse all incidences simultaneously for the negative face.
Each boundary vertex has one incoming and one outgoing edge,
so both incidence laws hold. The unoriented boundary is a closed
interval. It cannot be the boundary circle of a regular two-cell.
This example has four unoriented cells and eight oriented cells.

\subsection{A finite obstruction with all vertex links connected}
\begin{counterexample}\label{ex:torus-55}
There is a signed incidence system with $55$ unoriented cells,
of maximal dimension three, such that every signed diamond holds,
every boundary of positive dimension is connected, and every
vertex link in its order complex is connected, but it is not
the face poset of a regular CW complex.
\end{counterexample}
\begin{proof}
Triangulate the plane by unit squares split along their increasing
diagonals, and quotient by $3\Z^2$. The vertex set is
$V=(\Z/3)^2$. With $u=(1,0)$ and $v=(0,1)$, the two positively
oriented triangles based at $w\in V$ are
\[
 (w,w+u,w+u+v),\qquad (w,w+u+v,w+v).
\]
Every triangle embeds in the quotient because its coordinate span
is smaller than three. The two triangle types have distinct
translation orbits, each of size nine. The three edge directions
$u,v,u+v$ similarly give three orbits of nine edges. Thus this
torus triangulation has $9$ vertices, $27$ edges and $18$ triangles.
Adjoin a formal greatest element $\Omega$ of dimension three to
its nonempty face poset, obtaining $55$ cells.

Orient the edges arbitrarily and the triangles by the displayed
orientations. Give the incidence of each triangle into $\Omega$
coefficient $+1$. Every edge lies in two triangles with opposite
induced signs. Below a triangle, the ordinary simplex boundary
has the same cancellation at every vertex. These give every
rank-two diamond, together with the opposite endpoint signs on
edges. Doubling each cell by its two orientations gives the
required signed incidence system.

The two-cell boundaries are circles and the boundary of $\Omega$
is the connected torus. The link of the vertex $\Omega$ in the
order complex is the subdivided torus, hence connected. At an
old cell $c$, its vertex link is
\[
 \Delta(P_{<c})*\Delta(P_{>c}),
\]
where $P$ now includes $\Omega$. The upper interval has maximum
$\Omega$, so its order complex is a nonempty cone. This makes
the link connected; for a minimal cell the lower complex is
empty and the link is just the upper cone.

Finally the strict lower complex of $\Omega$ is a torus and has
first rational Betti number two, while a two-sphere has first
Betti number zero. Thus it cannot be the boundary of a regular
three-cell. The accompanying exact finite calculation gives
boundary ranks $8$ and $17$, hence Betti numbers $(1,2,1)$,
and verifies all $55$ vertex-link connectedness tests.
\end{proof}
The global links in this example must be distinguished from
the lower intervals appearing in Theorem~\ref{thm:cw}.


\section{Simplicial recognition and the two topologies}
\label{sec:simplicial-comparison}

\subsection{Connected boundaries}
For any poset $P$ and any $x\in P$, put $B_x=P_{<x}$.
The following elementary equivalence specifies the connectedness
condition that is weaker than the spherical hypothesis in
Theorem~\ref{thm:cw}.

\begin{proposition}\label{prop:boundary-connected}
The following are equivalent: $B_x$ is nonempty and connected
in either Alexandrov topology; its undirected comparability graph
is nonempty and connected; and $|\Delta(B_x)|$ is nonempty and
connected. If the principal ideal of $x$ is finite, the
comparability graph can be replaced by the Hasse graph.
\end{proposition}
\begin{proof}
A comparability component is both an upper and a lower set, hence
is clopen in either Alexandrov topology. Conversely a clopen
separation cannot separate comparable points: lower openness
of the two sides, or dually upper openness, rules this out.
The comparability edges are exactly the edges of the order
complex. Their graph components partition its simplices and give
clopen subcomplexes in the weak topology. A finite edge path is
a topological path, proving the second equivalence. In a finite
principal ideal, comparable points can be joined by a saturated
chain. The nonempty qualification excludes the empty boundary of
a vertex; the boundary of a regular edge is the disconnected
space $S^0$.
\end{proof}
Thus a connected-boundary axiom in dimensions at least two means
that any two proper faces are joined by a finite incidence zigzag
among proper faces. It is not a replacement for sphericality.

\subsection{The global intersection condition for simplices}
We prove the first assertion of
Theorem~\ref{thm:simplicial-comparison}. For the nonempty face
poset of an abstract simplicial complex, adjoining the empty face
gives binary meets by intersection and gives a finite Boolean
interval below each face.

Conversely assume the finite Boolean interval condition, and
write $V_x$ for the atoms of $P^0$ below $x$. If binary meets exist
and $V_x=V_y$, all these atoms lie below $x\wedge y$.
Within either Boolean interval they generate its top, so
$x=x\wedge y=y$. Thus the atom-set map is injective.
If instead injectivity is assumed, every subset of $V_x$ is
the atom set of an element of $[0,x]$. Hence the nonempty atom
sets form an abstract simplicial complex. Moreover
$V_x\subseteq V_y$ implies $x\le y$: the element below $y$
with atom set $V_x$ must be $x$. The atom-set map is an order
isomorphism to the nonempty face poset of that complex.
Intersections supply its binary meets. This proves all three
equivalent formulations, including the empty case.

Local Booleanity alone is insufficient. Two distinct edges with
the same two endpoints form a regular circle with Boolean
principal augmented intervals, but the two edge cells have the
same vertex set. Their common lower faces have no greatest
element, so adjoining zero does not give a meet-semilattice.

\begin{proposition}[A weaker predecessor condition]\label{prop:predecessor-repair}
The conjunction of binary meets in $P^0$ and the condition that
every nonminimal cell has at least two maximal proper faces does
not characterize simplicial complexes, even among finite regular
CW face posets. Replacing the predecessor condition by finite
Boolean principal intervals gives the exact criterion.
Equivalently, one may require finite distributive principal
intervals together with that predecessor condition.
\end{proposition}
\begin{proof}
Every positive-dimensional regular cell satisfies the predecessor
condition: a unique maximal proper face would make its boundary
order complex a cone, whereas that complex is a nonempty sphere.
The face poset of a square has binary intersections of closed
faces and satisfies the condition. Its augmented top interval
has ten elements and four atoms, so it cannot be the sixteen-element
Boolean lattice on those atoms. Even regularity does not follow
from the two weaker conditions: three incomparable vertices below
a single top cell satisfy them, but that cell has three boundary
points rather than the two of a regular one-cell.

In a finite distributive lattice every nonzero join-irreducible
has exactly one lower cover. The predecessor condition therefore
makes every such element an atom. Every element is a join of
join-irreducibles, so the lattice is atomistic. Distinct atoms have
zero meet; distributivity shows that different subsets of atoms
have different joins. Every subset has a join and every element
is such a join. This gives a Boolean lattice. The converse is
immediate, proving the final equivalence.
\end{proof}
In a regular cell complex, the meet condition says that two closed
cells intersect in a closed cell or not at all. It makes no claim
that the cells can be realized by convex polytopes.

\subsection{Flags and interval incidences}
The finite nonempty chains of a poset $L$, ordered by inclusion,
are exactly the nonempty faces of $\Delta(L)$. Subchains are
chains, and intersection gives the augmented meets. This is a
simplicial construction in every dimension.

The point-and-interval construction has a different incidence
relation. Its vertex cells are $P_a$ and its one-cells are
$I_{ab}$ for $a<b$, with the two endpoint incidences
$P_a,P_b<_{\mathrm{face}}I_{ab}$. It is the one-dimensional
simplicial complex of the comparability graph of $L$. It has
no higher-dimensional cells even when $L$ has long chains.
Its positional order, used in Theorem~\ref{thm:flags}, instead
puts $P_a<_{\mathrm p}I_{ab}<_{\mathrm p}P_b$.

A finite \emph{prefigure} in this construction is a positional
chain; a \emph{figure} is one closed under endpoint incidence.
The endpoint-closure operation is
\[
 F\longmapsto F\cup\bigcup_{I_{ab}\in F}\{P_a,P_b\}.
\]
A symbol preceding $I_{ab}$ has its last endpoint at most $a$,
and a symbol following it has its first endpoint at least $b$.
Thus endpoint closure preserves the positional-chain condition.
For example, $\{I_{ab}\}$ is a prefigure but is not a figure.
This also explains why the finite admissibility theorem was
stated for all pairwise orthogonal subsets, without assuming
incidence closure.

\subsection{The weak comparison with an Alexandrov space}
We finish Theorem~\ref{thm:simplicial-comparison} by checking the
local criterion in McCord's Theorem~6 \cite[p.~467]{McCord}.
That theorem says that a map is a weak equivalence if its
restrictions over a basis-like open cover are weak equivalences.
For $U_x=L_{\le x}$, the inverse image under $\mu_L$ consists
of points with a positive coordinate at some vertex of $U_x$.
It is open, being the complement of the full subcomplex on
$L\setminus U_x$. Renormalize the coordinates in $U_x$ and
linearly suppress the others. This deforms the inverse image
onto $|\Delta(U_x)|$, a cone with apex $x$. The same greatest
element contracts the Alexandrov space $U_x$. The principal
lower sets form a basis, so McCord's local theorem applies.
The coordinate deformation is continuous simplex by simplex
in the weak topology; no local finiteness is required.

This is generally only a weak equivalence. Take two incomparable
minima below two incomparable maxima. Its order complex is a
circle. Every continuous map from that finite connected poset
to a Hausdorff space is constant: comparable points must have
equal images. A homotopy inverse into the circle would therefore
make its identity null-homotopic, which is impossible.
For a regular CW face poset, the order complex instead gives
a barycentric subdivision homeomorphic to the original
Hausdorff realization. These comparisons have different targets.


\section{Scope}
The finite conditional-lattice theorem is not asserted for arbitrary infinite posets; the split-poset comparison and Boolean criterion are unrestricted. The powerset calculation uses primitive disjunction of open cells and must not be transferred without proof to a different whole-figure disjunction. For regular cell structures, spherical lower intervals concern cell boundaries, not global links: two triangles meeting only at a vertex already give a regular CW complex with a disconnected positive-dimensional global link at that vertex. Published-statement corrections are recorded separately from this paper.

\section*{Method and records}
The mathematics and text were produced by GPT-6 Astra (ChatGPT Pro mode). Checks are recorded with their actual scope and provenance; no human has checked the work. The text is dedicated to the public domain under CC0~1.0. The proposed repository for the accompanying programs and records is \url{\workbenchurl}; the workbench registry is \url{\recordurl}.
\begingroup\small
\begin{thebibliography}{99}
\bibitem{Groth} A. Grothendieck, \emph{Vers une g\'eom\'etrie des formes}, Montpellier archive, cote~156, especially parts~4, 7, 8 and 9; part~9, interval construction on archive sheets~4--7 (PDF pp.~5--8), and support axioms on archive sheets~50--52 (PDF pp.~51--53, manuscript pp.~47--49), spatial carriers and their target on archive sheets~45--49 (PDF pp.~46--50), and the ten-element reconstruction on archive sheet~112 (PDF p.~113). \url{https://grothendieck.umontpellier.fr/156-9.pdf}.
\bibitem{Groth7} A. Grothendieck, \emph{Vers une g\'eom\'etrie des formes}, Montpellier archive, cote~156, part~7, archive sheet~86 (PDF p.~87) and sheets~92--94 (PDF pp.~93--95), open supports and their proposed relative-complement calculus. \url{https://grothendieck.umontpellier.fr/156-7.pdf}.
\bibitem{Groth8} A. Grothendieck, \emph{Vers une g\'eom\'etrie des formes}, Montpellier archive, cote~156, part~8, archive sheets~84--85 (PDF pp.~85--86), definition of full faithfulness. \url{https://grothendieck.umontpellier.fr/156-8.pdf}.
\bibitem{Groth6} A. Grothendieck, \emph{Vers une g\'eom\'etrie des formes}, Montpellier archive, cote~156, part~6, archive and manuscript sheets~28--34 (PDF pp.~29--35); principal-multistrate membership on sheet~31, Th\'eor\`eme-Scholie on sheets~32--33. \url{https://grothendieck.umontpellier.fr/156-6.pdf}.
\bibitem{Hua} M. Hua, \emph{Grothendieck Montpellier. \'Enonc\'es des lectures modernis\'ees d\'emontr\'es}, Hexagon 2611.00125v2, 5 November 2026. Counterexample~79.1, pp.~369--370. \url{https://hexagonmath.org/2611.00125v2}.
\bibitem{Jenca} G. Jen\v ca, \emph{Orthogonality spaces associated with posets}, Order 40 (2023), 575--588. Theorem~4.10 in arXiv:2206.08113v1, p.~10. \url{https://doi.org/10.1007/s11083-022-09623-5}.
\bibitem{Jenca2026} G. Jen\v ca, \emph{N-free posets and orthomodularity}, Order 43 (2026), article 11. Theorem~3.2, p.~5; Theorem~4.2, p.~6. \url{https://doi.org/10.1007/s11083-025-09723-y}.
\bibitem{JencaCorrection} G. Jen\v ca, \emph{Correction: N-free posets and orthomodularity}, Order 43 (2026), article 28. Theorem~1; author version arXiv:2607.19423v1, pp.~2--4. \url{https://doi.org/10.1007/s11083-026-09742-3}; \url{https://arxiv.org/abs/2607.19423}.
\bibitem{Coste} M. Coste, \emph{An Introduction to O-minimal Geometry}, 1999. Theorems~3.4, 3.9, 3.22 and 4.4. \url{https://perso.univ-rennes1.fr/michel.coste/polyens/OMIN.pdf}.
\bibitem{Bjorner} A. Bj\"orner, \emph{Posets, regular CW complexes and Bruhat order}, European J. Combin. 5 (1984), 7--16, Definition~2.1 and Section~2.3, pp.~8--9; Proposition~3.1, p.~11. \url{https://doi.org/10.1016/S0195-6698(84)80012-8}.
\bibitem{NPW} M. Nielsen, G. Plotkin and G. Winskel, \emph{Petri nets, event structures and domains, part I}, Theoret. Comput. Sci. 13 (1981), 85--108, Definition~14, p.~99; Definitions~15--16, p.~101; Definitions~17--18 and Theorems~8--9, p.~102. \url{https://www.cl.cam.ac.uk/~gw104/PetriNetsPart1.pdf}.
\bibitem{Sumner} D. P. Sumner, \emph{Dacey graphs}, J. Austral. Math. Soc. \textbf{18} (1974), 492--502. Definition, p.~492. \url{https://doi.org/10.1017/S1446788700029232}.
\bibitem{PV} J. Paseka and T. Vetterlein, \emph{Categories of orthogonality spaces}, J. Pure Appl. Algebra \textbf{226} (2022), article~106859. Author version arXiv:2003.03313v2, Definition~2.4 and Lemma~2.5, pp.~5--6; Lemma~2.11 and Proposition~2.14, p.~8. \url{https://arxiv.org/abs/2003.03313v2}.
\bibitem{McCord} M. C. McCord, \emph{Singular homology groups and homotopy groups of finite topological spaces}, Duke Math. J. \textbf{33} (1966), 465--474. Theorems~2 and~6, pp.~466--467; Section~5, p.~469. \url{https://doi.org/10.1215/S0012-7094-66-03352-7}.
\end{thebibliography}
\endgroup
\end{document}
